\documentclass[11pt]{article}
\usepackage[a4paper,margin=2.5cm]{geometry}
\usepackage{amsmath,amssymb,amsthm}
\usepackage{longtable,array}
\usepackage[hidelinks,hypertexnames=false]{hyperref}

\newtheorem{theorem}{Theorem}
\newtheorem{lemma}[theorem]{Lemma}
\newtheorem{proposition}[theorem]{Proposition}
\newtheorem{corollary}[theorem]{Corollary}
\newtheorem{fact}{Fact}
\theoremstyle{remark}
\newtheorem*{remark}{Remark}

\emergencystretch=3em
\newcommand{\jp}{\circ}
\newcommand{\Pp}{P_{1/2}}
\newcommand{\opc}{\mathrel{|}}
\newcommand{\cA}{\mathcal A}
\newcommand{\cP}{\mathcal P}\newcommand{\cQ}{\mathcal Q}\newcommand{\cL}{\mathcal L}
\newcommand{\cR}{\mathcal R}\newcommand{\cK}{\mathcal K}
\DeclareMathOperator{\Exp}{Exp}
\DeclareMathOperator{\spec}{sp}
\DeclareMathOperator{\comm}{comm}
% Provenance annotations: \prov{status: where}{one-line justification}
\newcommand{\prov}[2]{\par\smallskip\noindent{\small\textbf{[#1]}\ #2}\par\smallskip}

\title{The commutation theorem for sequential products\\ in JB-algebras: a detailed proof}
\author{}
\date{26 September 2026}

\begin{document}
\maketitle

\begin{abstract}
We give a complete proof, at the level of individual identities and
inequalities, that $U_{\sqrt a}\,b = U_{\sqrt b}\,a$ with $a,b\ge 0$ forces $b$ to
operator-commute with the closed subalgebra $C(a)$, first in JBW-algebras and then in
arbitrary JB-algebras. The argument follows the machine-checked proof step by step.
No trace, no Shirshov--Cohn theorem and no structure theory are used. Macdonald's theorem
is avoided too: the two standard facts that the textbook proofs derive from it, the
fundamental formula and the positivity of $U_a$, are proved in \hyperref[app:mac]{Appendix~A}
from one identity and analysis. Beyond the standard facts of \S2, all algebra comes from
the linearised Jordan identity.
\end{abstract}

\section{Statement and notation}

For $a,b$ in a JB-algebra $\cA$ write $L_a b = a\jp b$,
\[
  U_a = 2L_a^2 - L_{a^2},\qquad U_{a,c} = L_aL_c + L_cL_a - L_{a\jp c},
\]
so $U_a = U_{a,a}$, $U_{a+c} = U_a + U_c + 2U_{a,c}$, $U_{1,c} = L_c$, and
$U_ab = aba$, $U_{a,c}b = \tfrac12(abc+cba)$ in a special algebra.
Elements $a,b$ \emph{operator-commute}, written $a\opc b$, if $L_aL_b = L_bL_a$.
For $a\in\cA$, $C(a)$ is the closed unital subalgebra generated by $a$, and
$\comm(a) = \{b : b\opc c \text{ for all } c\in C(a)\}$.

\begin{theorem}\label{thm:main}
Let $a,b\geq 0$ in a JB-algebra $\cA$. If $U_{\sqrt a}\, b = U_{\sqrt b}\, a$, then
$b\in\comm(a)$; in particular $a\opc b$.
\end{theorem}
\prov{Their work: vdW 2020, Thms.~3.12--3.13}{The statement is van de Wetering's
\cite{vdW}, proved there via Shirshov--Cohn \cite[7.2.5]{HOS}, Fuglede--Putnam--Rosenblum
\cite{Fug50,Put51,Ros58},
Shultz's structure theorem \cite{Shu79} and the bidual; for effects on a Hilbert space it
is Gudder--Nagy \cite{GN01}. The proof below is different; see \S\ref{sec:prov}.}

This is the commutation property needed for $[0,1]_\cA$ with $a\ast b = U_{\sqrt a}b$
to be a sequential effect algebra (van de Wetering \cite{vdW}; see the Remark at the end of
Part~VI). Conversely, if $b\in\comm(a)$ then $b\opc\sqrt a$ (as $\sqrt a\in C(a)$), so
$U_{\sqrt a}b = a\jp b$ by Fact~\ref{f:Uopc}. Moreover $b\opc a$, so the powers argument of
VI.8 (with $b$ in the role of $a$ and $a$ in the role of $y$) gives $b^n\opc a$ for all $n$;
as $\sqrt b\in C(b)$ is a norm limit of polynomials in $b$ and $\{c : c\opc a\}$ is closed,
$\sqrt b\opc a$, and $U_{\sqrt b}a = b\jp a$ by Fact~\ref{f:Uopc} again. So
$U_{\sqrt a}b = U_{\sqrt b}a$.
Throughout, $x = \sqrt a$ and $y = \sqrt b$, so that the hypothesis reads
$U_x(y^2) = U_y(x^2)$, and
\[
  A = x\jp y,\qquad D = [L_x,L_y].
\]

\paragraph{Notation.}
Brackets are commutators of operators on $\cA$, and an operator applied
to an element is written $T(z)$ or $Tz$.
We use two different exponentials.
For an element $h$, $e^{h}\in C(h)$ is the functional-calculus exponential;
for a bounded operator $T$ on $\cA$, $\Exp(T) = \sum_n T^n/n!$ is the operator
exponential.
They are related by $U_{e^{h}} = \Exp(2L_h)$.
For $h\in\cA$ and $f\in C(\spec h)$ we write $f(h)$ for the functional calculus.
All operators below are bounded, and $B(\cA)$ is the Banach algebra of bounded
operators with the operator norm.

\subsection{Provenance of the ingredients}\label{sec:prov}
Each step below carries a bracketed tag: \textbf{[Their work]} (a paper or thesis of
A.~Westerbaan, B.~Westerbaan or J.~van de Wetering), \textbf{[Literature]}, or
\textbf{[New?]}. \textbf{[New?]} only means that \emph{we have not found it}; it is not a
priority claim. What we searched: Hanche-Olsen--St{\o}rmer \cite{HOS}; McCrimmon's
\emph{Taste} \cite{McC} (in particular the Fundamental Formulas 5.2.3 and the
commutator calculus of Ch.~III.22); Alfsen--Shultz \cite{AS}; van de Wetering
\cite{vdW,vdW19,vdWFF}; the follow-up papers \cite{WWvdW20,WWvdW21}; the two
Westerbaan theses (which treat the sequential product only in von Neumann algebras, and
contain no Jordan-algebraic ingredient used here); and web searches for a Jordan
Fuglede--Putnam theorem (``Fuglede--Putnam'' with ``Jordan algebra'', ``JB-algebra'',
``JB$^*$-triple''), which found nothing beyond \cite{vdW}. HOS section numbers were
checked against the authors' online edition; numbers marked ``unverified'' were not.
\begingroup\small
\renewcommand{\arraystretch}{1.15}
\begin{longtable}{>{\raggedright}p{4.1cm}>{\raggedright}p{7.6cm}>{\raggedright\arraybackslash}p{2.9cm}}
\hline
\textbf{Ingredient} & \textbf{Where} & \textbf{Status}\\\hline\endhead
Theorem~\ref{thm:main} (statement) & vdW \cite[Thms.~3.12, 3.13]{vdW}; Hilbert-space
effects: Gudder--Nagy \cite{GN01}; JC-algebras: \cite[Prop.~3.3]{vdW} & Their work\\
SEA application & vdW \cite[Thm.~4.2, Prop.~4.4]{vdW}; axioms of Gudder--Greechie
\cite{GG02}; used in \cite{WWvdW20,WWvdW21} & Their work\\
Facts \ref{f:norm}, \ref{f:cfc}, \ref{f:jbw} (norm, functional calculus, JBW) &
\cite[\S\S3.1, 3.2, 4.1, 4.2]{HOS}; constant $3$: \cite[3.2.7]{HOS}; idempotents dense:
\cite[4.2.3]{HOS}. The form of $\pi_\rho$ and \eqref{eq:specapprox} is a routine
rewriting, not stated so in HOS; proof sketches after Fact~\ref{f:jbw} & Literature\\
Fact \ref{f:pos} (positivity of $U_a$) & \cite[3.3.6]{HOS} (proved there via 2.4.17,
i.e.\ Macdonald); finite-dimensional analytic proof via $P(e^x) = e^{2L(x)}$:
Faraut--Kor\'anyi \cite[Ch.~II--III, numbers unverified]{FK}. Macdonald-free proof here:
\hyperref[app:mac]{Appendix~A}, A.5--A.6 (resolvent positivity of $L_h$, cf.\ Alfsen--Shultz's
criterion \cite[(1.82), number unverified]{ASss}; then the swap identity) & Literature;
\hyperref[app:mac]{App.~A} route New?\\
Fact \ref{f:ff} (fundamental formula) & \cite[2.4.18]{HOS} via Macdonald \cite{Mac60};
Macdonald-free: vdW \cite{vdWFF} (algebraic), \cite{FK} (analytic, finite-dim.); here
\hyperref[app:mac]{Appendix~A}, A.1--A.4 &
Literature / their work; \hyperref[app:mac]{App.~A} route New?\\
$U_{e^h} = \Exp(2L_h)$ and the ODE proof of FF (\hyperref[app:mac]{Appendix~A}) &
Identity: \cite{FK} for Euclidean Jordan algebras; $2U_{w,k\jp w} = L_kU_w+U_wL_k$ is
McCrimmon's (FFV) at $y=1$ \cite[5.2.3]{McC}. The Banach-space ODE argument for
JB-algebras: not found & Literature; ODE route New?\\
Fact \ref{f:S3} (orthogonality) & Alfsen--Shultz \cite[Lemma~1.26, as quoted in
{\cite[Lemma~4.1]{vdW}}]{AS} & Literature\\
Fact \ref{f:peirce} (Peirce) & \cite[\S2.6]{HOS}; $\{p\}'$ a subalgebra: Jacobson
\cite{Jac89}, \cite[Prop.~2.13]{vdW} & Literature\\
Lemma \ref{lem:peirce} ($U_pU_z(1-p) = 0\Rightarrow\Pp z = 0$) & elementary Peirce
calculus; not located as a stated lemma & New? (minor)\\
(J), (II) & \cite[\S2.4]{HOS}, \cite{Jac68}; \cite[Lemma~2.4(c), Prop.~2.5]{vdW} &
Literature / their work\\
Lemma \ref{lem:der} ($[L_x,L_y]$ derivation) & Jacobson \cite{Jac49}; \cite[Problem
II.3.6(3), Ch.~III.22.1]{McC} & Literature\\
Lemma \ref{lem:1} ($D(A) = \frac14(U_xy^2-U_yx^2)$) & associative shadow ($xy$ normal
iff $xy^2x = yx^2y$) used in \cite[Props.~3.2, 3.3]{vdW}; the Jordan identity not found
(McCrimmon's ``commutator square'' \cite[22.1.1]{McC} is a different combination) &
New? (as a Jordan identity)\\
Lemma \ref{lem:3} ($2[L_x,L_{x\jp y}] + [L_y,L_{x^2}] = 0$) & \cite[Lemma~2.4(b)]{vdW};
standard linearisation & Their work / literature\\
Lemma \ref{lem:2} (intertwining) & is McCrimmon's Commuting Formula (FFII)
$V_{x,y}U_x = U_xV_{y,x}$ \cite[5.2.3]{McC}, a quadratic Jordan axiom
\cite{McC66,Jac69}, since $V_{x,y} = 2(L_{x\jp y}+[L_x,L_y])$. Only the hand derivation
from four instances of (J) is ours & Literature\\
Lemma \ref{lem:expder} ($\Exp$ of a derivation) & standard (automorphism: any Banach
algebra; contractivity: unital positive maps of order-unit spaces) & Literature\\
Prop.~\ref{prop:flow} (the flow) & Jordan analogue of Rosenblum's exponential proof of
Fuglede--Putnam \cite{Ros58}; no Jordan version found & New?\\
Lemma \ref{lem:gap}, Props.~\ref{prop:stepB}, \ref{prop:cp}, Lemma \ref{lem:twosided} &
the idea (a bounded two-sided exponential orbit kills off-diagonal spectral blocks) is
standard for operators; the Jordan form with $U$-compressions not found & New?\\
Prop.~\ref{prop:pows} & spectral approximation, cf.\ \cite[4.2.3]{HOS},
\cite[Props.~2.33--2.35]{vdW} & Literature\\
Prop.~\ref{prop:ya} (Jordan Fuglede--Putnam, JBW) & statement follows from
\cite[Thm.~3.12]{vdW}; proof without Shirshov--Cohn/structure theory not found &
New? (proof)\\
Prop.~\ref{prop:trick} ($(y+\lambda,a)$ trick) & --- & New? (minor)\\
Prop.~\ref{prop:mulmem} & \cite[Prop.~2.34]{vdW} (same route: idempotents and
\cite{Jac89}) & Their work\\
Lemma \ref{lem:SQ} ((SQ)) & statement: \cite[Thm.~3.13, (a)$\Leftrightarrow$(d)]{vdW},
via structure theory (cf.\ his Remark~2.15); this proof not found & Their work
(statement); New? (proof)\\
VI.1 ($U$ multiplicative on $C(h)$) & consequence of the fundamental formula on an
associative subalgebra & Literature (routine)\\
VI.2 ($U_ce = 0\Rightarrow c\jp e = 0$) & for $c\ge0$ in \cite[Lemma~1.26]{AS}; the
$U_{1+tc}$ argument is routine & Literature (variant)\\
VI.3 (cross term) & linearisation of FF applied to $1$ & Literature (routine)\\
VI.4, VI.5 (off-diagonal parts, second-order vanishing) & --- & New?\\
VI.6 (Kleinecke--Shirokov) & Kleinecke \cite{Kle57}, Shirokov \cite{Shi56}; the bound
is their standard proof & Literature\\
VI.7 (Gelfand--Hille) & Gelfand \cite{Gel41}, Hille \cite{Hil44}, survey Zem\'anek
\cite{Zem94}; here a special case with Phragm\'en--Lindel\"of & Literature\\
VI.8 and the JB proof & result: \cite[Thm.~3.13]{vdW} via $\cA^{**}$; the bidual-free
route not found & New? (proof)\\
Final Remark ($\comm(c)$ a subalgebra) & \cite[Prop.~2.34]{vdW} via $\cA^{**}$ & Their
work (statement); New? (proof)\\\hline
\end{longtable}
\endgroup
\noindent Note: as proved in \cite{HOS}, Facts~\ref{f:pos} and~\ref{f:ff} rely on
Macdonald's theorem; \hyperref[app:mac]{Appendix~A} gives Macdonald-free proofs of both, following
the formalisation.

\section{Standard facts}

We use the following facts; references are to Hanche-Olsen and St{\o}rmer
\cite{HOS}.

\begin{fact}[order-unit norm, {\cite[\S3.1]{HOS}}]\label{f:norm}
$-\|v\|1\le v\le \|v\|1$; if $-\lambda 1\le v\le\lambda 1$ then $\|v\|\le\lambda$;
hence $0\le v\le w$ implies $\|v\|\le\|w\|$. Moreover $\|a\jp b\|\le\|a\|\|b\|$,
$\|a^2\| = \|a\|^2$, and squares are positive.
Consequently $\|U_cv\|\le 3\|c\|^2\|v\|$ and $\|U_{a,c}w\|\le 3\|a\|\|c\|\|w\|$
(triangle inequality on the three terms; cf.\ \cite[3.2.7]{HOS}).
\end{fact}
\prov{Literature: HOS \S3.1, 3.2.7}{Standard order-unit and JB-norm facts; the constant $3$ is \cite[3.2.7]{HOS}.}

\begin{fact}[functional calculus, {\cite[\S3.2]{HOS}}]\label{f:cfc}
$C(h)$ is associative, $\spec h\subseteq[-\|h\|,\|h\|]$, and $f\mapsto f(h)$ is an
isometric isomorphism $C(\spec h)\to C(h)$ with $f(h)\ge0$ iff $f\ge0$ on $\spec h$
(so $f\le g$ iff $f(h)\le g(h)$). Any two elements of $C(h)$ operator-commute.
Jordan powers are associative: $h^i\jp h^j = h^{i+j}$.
For $c\in\mathbb R$, $e^{ch} := \sum_n c^nh^n/n! = \Exp(cL_h)(1)$ equals
$(t\mapsto e^{ct})(h)$. Every $a\ge0$ has a unique square root $\sqrt a\ge 0$, and it lies in $C(a)$.
\end{fact}
\prov{Literature: HOS \S3.2}{Continuous functional calculus; power-associativity also \cite[Cor.~2.9]{vdW}.}

\begin{fact}\label{f:Uopc}
If $r\opc v$ then $U_rv = r^2\jp v$. (Apply $L_rL_v = L_vL_r$ to $r$:
$r\jp(r\jp v) = r\jp(v\jp r) = v\jp r^2$, so $U_rv = 2r^2\jp v - r^2\jp v$.)
\end{fact}
\prov{Elementary}{One line from the definition.}

\begin{fact}[positivity, {\cite[3.3.6]{HOS}}]\label{f:pos}
$U_a v\ge 0$ whenever $v\ge 0$, for every $a\in\cA$. Hence each $U_a$ is monotone.
\end{fact}
\noindent A Macdonald-free proof is given in \hyperref[app:mac]{Appendix~A} (A.5--A.6).
\prov{Literature: HOS 3.3.6}{Proved in \cite{HOS} via 2.4.17, which rests on Macdonald's theorem; for Euclidean Jordan algebras an analytic proof via $P(e^x) = e^{2L(x)}$ is in \cite{FK} (numbers unverified). The proof in \hyperref[app:mac]{Appendix~A} avoids Macdonald.}

\begin{fact}[fundamental formula, {\cite[\S2.4]{HOS}}]\label{f:ff}
$U_{U_ab} = U_aU_bU_a$; with $b=1$, $U_{a^2} = U_a^2$. Moreover
$U_{e^h} = \Exp(L_h)^2 = \Exp(2L_h)$.
\end{fact}
\prov{Literature: HOS 2.4.18; their work: vdW 2018}{FF via Macdonald \cite{Mac60} in \cite{HOS}; a Macdonald-free algebraic proof is \cite{vdWFF}. $U_{e^h} = \Exp(2L_h)$ is not in HOS \S2.4; for Euclidean Jordan algebras it is in \cite[Ch.~II]{FK}.}

\begin{remark}
Both parts of Fact~\ref{f:ff}, and Fact~\ref{f:pos}, are proved without Macdonald's theorem
in \hyperref[app:mac]{Appendix~A}, following the formalisation: two instances of the
identity~(J) below give $2U_{w,k\jp w} = L_kU_w + U_wL_k$; hence
$t\mapsto \Exp(-tL_h)U_{\Exp(tL_h)b}\Exp(-tL_h)$ has zero derivative, which gives
$U_{\Exp(L_h)b} = \Exp(L_h)U_b\Exp(L_h)$ and, at $b = 1$, $U_{e^h} = \Exp(2L_h)$. The
fundamental formula then holds at every $a = e^{h}$, i.e.\ at every $a\ge\delta 1$
with $\delta > 0$, and it extends to all $a$ because $\lambda\mapsto U_{U_{a+\lambda}b} -
U_{a+\lambda}U_bU_{a+\lambda}$ is a polynomial that vanishes for $\lambda>\|a\|$.
\end{remark}
\prov{New? (route)}{The identity used is McCrimmon's (FFV) at $y=1$ \cite[5.2.3]{McC}; analytic proofs of FF exist in finite dimension \cite{FK}; we have not found this Banach-space ODE argument for JB-algebras.}

\begin{fact}[orthogonality of $U$]\label{f:S3}
If $a,b\ge 0$ and $U_{\sqrt a}b = 0$, then $U_{\sqrt b}a = 0$.
\end{fact}
\prov{Literature: Alfsen--Shultz Lemma 1.26}{As quoted in \cite[Lemma~4.1]{vdW}; the proof below via FF is standard.}
\begin{proof}
Put $r=\sqrt a$, $s=\sqrt b$. By Fact~\ref{f:ff},
$(U_s r^2)^2 = U_{U_sr^2}1 = U_sU_{r^2}U_s 1 = U_sU_rU_r(s^2) = U_sU_r(U_rb)=0$.
So $\|U_sr^2\|^2 = \|(U_sr^2)^2\| = 0$.
\end{proof}

\begin{fact}[Peirce decomposition, {\cite[\S2.6]{HOS}}]\label{f:peirce}
Let $p$ be an idempotent. Then $P_1 = U_p = 2L_p^2-L_p$, $\Pp = 4L_p - 4L_p^2$ and
$P_0 = U_{1-p}$ are complementary projections onto $\cA_i = \{z: p\jp z = iz\}$,
with $\cA_1\jp\cA_1\subseteq \cA_1$, $\cA_0\jp\cA_0\subseteq\cA_0$,
$\cA_1\jp\cA_0 = 0$, $\cA_{1/2}\jp\cA_{1/2}\subseteq\cA_1+\cA_0$,
$\cA_i\jp\cA_{1/2}\subseteq \cA_{1/2}$ ($i = 0,1$).
Hence $\ker\Pp(p) = \cA_1+\cA_0$ is a subalgebra, and $\Pp(p)z = 0$ implies $p\opc z$.
Conversely $p\opc z$ implies $\Pp(p)z=0$: applying $L_pL_z = L_zL_p$ to $p$ gives
$p\jp(p\jp z) = p\jp z$. Finally, for idempotents, $p\le q$ iff $q\jp p = p$.
\end{fact}
\prov{Literature: HOS \S2.6}{Peirce rules; $\ker\Pp(p)$ a subalgebra is Jacobson's lemma \cite{Jac89}, cf.\ \cite[Prop.~2.13]{vdW}.}

\begin{lemma}\label{lem:peirce}
If $p$ is an idempotent and $U_pU_z(1-p) = 0$, then $\Pp(p)z = 0$.
\end{lemma}
\prov{New? (minor)}{Elementary Peirce calculus; we did not find it stated as a lemma (not in HOS \S2.6).}
\begin{proof}
Write $z = z_1+z_{1/2}+z_0$ with $z_i = P_i(p)z$. Since $U_p = P_1$, we have $U_pz = z_1$,
and the fundamental formula applied to $1$ gives
$U_pU_zp = U_pU_zU_p1 = U_{U_pz}1 = z_1^2$. By the Peirce rules
$P_1(z^2) = z_1^2 + P_1(z_{1/2}^2)$ (the products $z_1\jp z_0$ vanish, and
$z_{1/2}\jp z_i$, $z_0^2$ have no $\cA_1$-part). Using $U_z(1-p) = z^2 - U_zp$,
\[
 0 = U_pU_z(1-p) = P_1(z^2) - z_1^2 = P_1(p)(w^2) = P_0(1-p)(w^2),\qquad w = z_{1/2},
\]
and $w\in\cA_{1/2}(1-p)$. Put $q = 1-p$. Then $w^2\in \cA_1(q)+\cA_0(q)$ and its
$\cA_0(q)$-part vanishes, so $q\jp w^2 = w^2$. The Jordan identity
$(w\jp q)\jp w^2 = w\jp(q\jp w^2)$ reads $\tfrac12 w\jp w^2 = w\jp w^2$, so
$w^3 = 0$, $w^4 = w\jp w^3 = 0$, and $\|w\|^4 = \|w^2\|^2 = \|w^4\| = 0$.
\end{proof}

\begin{fact}[JBW-algebras, {\cite[\S4.1--4.2]{HOS}}]\label{f:jbw}
Let $\cA$ be a JBW-algebra. Bounded increasing sequences have suprema, and
(\emph{normality}) for every $r\in\cA$ the map $U_r$ preserves the supremum of a bounded
increasing sequence in $[0,1]$; if $z$ operator-commutes with every member of such a
sequence it operator-commutes with its supremum. For $h\in\cA$ and $\rho\in C(\spec h)$,
$\rho\ge0$, the element
\[
  \pi_\rho := \sup_n\, \sigma_n(h),\qquad \sigma_n = \min((n+1)\rho,1),
\]
is an idempotent (the spectral projection $\chi_{\{\rho>0\}}(h)$).
For $\delta>0$ and $j\in\mathbb N$ let $\rho_j(t) = (t+\|h\|-j\delta)^+$ and
$p_j = \pi_{\rho_j}$. Then $p_j\le p_i$ for $i\le j$, and for some $m$ and
$c\in\mathbb R$,
\begin{equation}\label{eq:specapprox}
  \Bigl\|h - \Bigl(c1 + \delta\sum_{i<m} p_{i+1}\Bigr)\Bigr\|\le\delta .
\end{equation}
\end{fact}
\prov{Literature: HOS \S4.1--4.2}{Monotone completeness, normality of $U_r$, idempotents dense \cite[4.2.3]{HOS}; the specific $\pi_\rho$ and \eqref{eq:specapprox} are routine rewritings not stated in this form in HOS.}

\paragraph{Proof sketches for Fact~\ref{f:jbw}.} We sketch the three parts that are not
stated in this form in \cite{HOS}, following the formalisation. Existence of suprema of
bounded increasing sequences, normal states, the fact that normal states separate points,
and normality of $U_r$ for general $r$ are taken from \cite[\S4.1]{HOS}; below only $r\ge0$
is needed, and for it (iii) gives a direct proof.

(i) \emph{$\pi_\rho$ is an idempotent.} Put $x_n = \sigma_n(h)$; these lie in $[0,1]$ and
increase, so $p := \pi_\rho = \sup_n x_n$ exists. Pointwise
$0\le\sigma_n-\sigma_n^2\le\sigma_{2n+1}-\sigma_n$: with $u = (n+1)\rho(t)$ both sides
vanish if $u\ge1$, and for $u<1$ we have $\sigma_{2n+1} = \min(2u,1)$ and
$u-u^2\le\min(2u,1)-u$. Let $\omega$ be a normal state; then $\omega(x_n)\to\omega(p)$, so
$\omega(x_n - x_n^2)\le\omega(x_{2n+1}) - \omega(x_n)\to0$. Also
$p^2 - x_n^2 = (p+x_n)\jp(p-x_n)$, and by the Cauchy--Schwarz inequality
$|\omega(u\jp v)|^2\le\omega(u^2)\omega(v^2)$, with $\omega((p+x_n)^2)\le4$ and
$\omega((p-x_n)^2)\le\omega(p-x_n)$ (as $0\le p-x_n\le1$), we get $\omega(x_n^2)\to\omega(p^2)$.
Hence $\omega(p^2) = \lim\omega(x_n^2) = \lim\omega(x_n) = \omega(p)$ for every normal state,
and $p^2 = p$. Since $\rho_j\le\rho_i$ for $i\le j$, $\sigma_n(\rho_j)(h)\le\sigma_n(\rho_i)(h)\le
p_i$ for all $n$, so $p_j\le p_i$.

(ii) \emph{The approximation \eqref{eq:specapprox}.} Let $\kappa_j(t) = \min(\rho_j(t),\delta)$.
Pointwise, $\delta\,\sigma_n(\rho_{j+1})\le\kappa_j$ (the left side is non-zero only where
$t+\|h\|-j\delta>\delta$, and there $\kappa_j = \delta$), and
$\kappa_j\le\delta\,\sigma_n(\rho_j)+\tfrac1{n+1}$ (if $\rho_j(t)\ge\frac1{n+1}$ then
$\sigma_n(\rho_j)(t) = 1$; otherwise $\kappa_j\le\rho_j<\frac1{n+1}$). Applying the
functional calculus, taking the supremum over $n$ in the first inequality and letting
$n\to\infty$ in the second,
\[
 \delta\,p_{j+1}\le\kappa_j(h)\le\delta\,p_j .
\]
Let $m = \lceil 2\|h\|/\delta\rceil$, so $(m+1)\delta>2\|h\|$. For real $u$ and $K\in\mathbb N$,
$\sum_{j<K}\min((u-j\delta)^+,\delta) = \min(u^+,K\delta)$; with $u = t+\|h\|\in[0,2\|h\|]$
and $K = m+1$ this gives $\sum_{j\le m}\kappa_j(h) = h+\|h\|1$. Hence, using
$\kappa_m(h)\ge0$ and $p_0\le1$,
\[
 \delta\sum_{i<m}p_{i+1}\le\sum_{j<m}\kappa_j(h)\le h+\|h\|1\le\delta\sum_{j\le m}p_j
 \le\delta1+\delta\sum_{i<m}p_{i+1},
\]
i.e.\ $0\le h-(c1+\delta\sum_{i<m}p_{i+1})\le\delta1$ with $c = -\|h\|$, and
Fact~\ref{f:norm} gives~\eqref{eq:specapprox}. (In $C(h)$-terms the approximant is the
function $-\|h\|+\delta\cdot\#\{k\ge1 : k\delta-\|h\|<t\}$, which lies in $[t-\delta,t]$.)

(iii) \emph{Suprema, in any JB-algebra.} Let $D\subseteq[0,1]$ be non-empty and directed
with supremum $s$. (a) \emph{$U_r$ preserves $\sup D$ for $r\ge0$.} By scaling
($U_{\lambda r} = \lambda^2U_r$) let $0\le r\le1$. For $\varepsilon>0$, $r+\varepsilon1\ge
\varepsilon1$, so by A.3 and A.2 of \hyperref[app:mac]{Appendix~A} $U_{r+\varepsilon1} = \Exp(2L_g)$
for some $g$, with inverse $\Exp(-2L_g)$; both are positive by A.5, so $U_{r+\varepsilon1}$
is an order automorphism and preserves suprema. Since $U_{r+\varepsilon1}v - U_rv =
2\varepsilon\,r\jp v+\varepsilon^2v$ has norm $\le\eta := 2\varepsilon+\varepsilon^2$ for
$v\in[0,1]$, an upper bound $u$ of $U_r(D)$ gives $U_{r+\varepsilon1}d\le u+\eta1$ for
$d\in D$, hence $U_{r+\varepsilon1}s\le u+\eta1$ and $U_rs\le u+2\eta1$; letting
$\varepsilon\to0$, $U_rs\le u$. So $U_rs = \sup U_r(D)$. (b) \emph{Operator commutation
passes to $s$.} Call a positive operator \emph{good} if it preserves suprema of non-empty
bounded directed subsets of $\cA_+$; sums, products and non-negative multiples of good
operators are good, and $U_r$ ($r\ge0$) is good by (a). For every $v$ and $N\ge\|v\|$ the
elements $\frac12(N1\pm v)$ are $\ge0$, and
\[
 N\,L_v = U_{\frac12(N1+v)} - U_{\frac12(N1-v)} ,
\]
so each $L_v$, and hence each $T = L_zL_w - L_{z\jp w}$, is a difference $P-P'$ of good
operators. If $z\opc d$ for all $d\in D$ then $T(d) = z\jp(w\jp d) - (z\jp w)\jp d = 0$ for
all $d\in D$, so $P(s) = \sup P(D) = \sup P'(D) = P'(s)$, i.e.\ $T(s) = 0$:
$z\jp(s\jp w) = s\jp(z\jp w)$ for every $w$, which is $z\opc s$. In particular $z$
operator-commutes with the supremum of a bounded increasing sequence in $[0,1]$ whenever it
operator-commutes with each member.
\prov{Literature (routine); formalised}{(i) is \texttt{sproj}, \texttt{sproj\_idem}; (ii) is
\texttt{clamp\_le\_proj}, \texttt{proj\_le\_clamp}, \texttt{sum\_clamp},
\texttt{spectral\_approx\_norm} in \texttt{JBWProj} and \texttt{spectral\_approx\_expl} in
\texttt{JBCommThm}; (iii) is \texttt{jQ\_isLUB} (\texttt{JBAllSEA}) and \texttt{nrm\_mulC},
\texttt{Nrm.eq\_zero}, \texttt{opComm\_of\_isLUB} (\texttt{JBCommThm}, \S1).}

Finally, for fixed $z$ the set $\{c : c\opc z\}$ is a closed linear subspace containing $1$
(the product is jointly continuous).

\section*{Part I. Identities}

\paragraph{The linearised Jordan identity.}
The Jordan identity $a^2\jp(a\jp w) = a\jp(a^2\jp w)$ says $[L_a,L_{a^2}] = 0$.
Replace $a$ by $\lambda a+\mu b+\nu c$: the left side of
$[L_{\lambda a+\mu b+\nu c}, L_{(\lambda a+\mu b+\nu c)^2}]=0$ is a polynomial in
$\lambda,\mu,\nu$ that vanishes identically, so its $\lambda\mu\nu$-coefficient
$2\bigl([L_a,L_{b\jp c}]+[L_b,L_{c\jp a}]+[L_c,L_{a\jp b}]\bigr)$ vanishes. In element
form, for all $a,b,c,d$:
\begin{equation}\tag{J}
  j(a,b,c;d) := \sum_{\mathrm{cyc}(a,b,c)} \Bigl(a\jp\bigl((b\jp c)\jp d\bigr) -
  (b\jp c)\jp(a\jp d)\Bigr) = 0 .
\end{equation}
Read as an operator in $d$ this is
$[L_a,L_{b\jp c}]+[L_b,L_{c\jp a}]+[L_c,L_{a\jp b}] = 0$; read as an operator in $a$,
it is
\begin{equation}\tag{II}
  L_{(b\jp c)\jp d} = L_{b\jp c}L_d + L_{b\jp d}L_c + L_{c\jp d}L_b
  - L_bL_dL_c - L_cL_dL_b .
\end{equation}
Every identity below is a rational linear combination of instances of~(J), with
commutativity used freely.
\prov{Literature; their work: vdW 2020, Lemma 2.4(c), Prop.~2.5}{(J) and (II) are the standard linearised Jordan identities \cite[\S2.4]{HOS}, \cite{Jac68}; (II) is \cite[Prop.~2.5]{vdW}.}

\begin{lemma}[$D$ is a derivation]\label{lem:der}
For all $x,y,u,v$, $D = [L_x,L_y]$ satisfies $D(u\jp v) = D(u)\jp v + u\jp D(v)$.
Consequently $D(1)=0$ and $[D,L_z] = L_{D(z)}$ for every $z$.
\end{lemma}
\prov{Literature: Jacobson 1949}{$[L_x,L_y]$ is the classical inner derivation \cite{Jac49}; cf.\ \cite[Problem II.3.6(3)]{McC}.}
\begin{proof}
Write out the two instances
\begin{align*}
 j(v,x,u;y) &= v\jp((x\jp u)\jp y) - (x\jp u)\jp(v\jp y) + x\jp((u\jp v)\jp y)
   - (u\jp v)\jp(x\jp y) \\&\quad + u\jp((v\jp x)\jp y) - (v\jp x)\jp(u\jp y),\\
 j(u,y,v;x) &= u\jp((y\jp v)\jp x) - (y\jp v)\jp(u\jp x) + y\jp((v\jp u)\jp x)
   - (v\jp u)\jp(y\jp x) \\&\quad + v\jp((u\jp y)\jp x) - (u\jp y)\jp(v\jp x).
\end{align*}
In $j(v,x,u;y) - j(u,y,v;x)$ the six products of two products cancel in pairs
($(x\jp u)\jp(v\jp y)$ against $(y\jp v)\jp(u\jp x)$, and so on), and the remaining six terms
are
\[
 x\jp(y\jp(u\jp v)) - y\jp(x\jp(u\jp v)) - v\jp\bigl(x\jp(y\jp u)-y\jp(x\jp u)\bigr)
 - u\jp\bigl(x\jp(y\jp v) - y\jp(x\jp v)\bigr),
\]
which is $D(u\jp v) - D(u)\jp v - u\jp D(v)$. By (J) it is $0$.
Taking $u=v=1$ gives $D(1) = 2D(1)$. Finally
$[D,L_z]w = D(z\jp w) - z\jp D(w) = D(z)\jp w$.
\end{proof}

\begin{lemma}[the value of $D$ at $A$]\label{lem:1}
For all $x,y$: $\;D(A) = [L_x,L_y](x\jp y) = \tfrac14\bigl(U_x(y^2) - U_y(x^2)\bigr)$.
\end{lemma}
\prov{New? (as a Jordan identity)}{In a special algebra it says $[[x,y],xy+yx] = 2(xy^2x-yx^2y)$, i.e.\ $xy$ is normal iff $xy^2x = yx^2y$, which is used in \cite[Props.~3.2, 3.3]{vdW}; we did not find the Jordan form (searched HOS \S2.4, McCrimmon Ch.~III.22).}
\begin{proof}
The two instances are
\begin{align*}
 j(y,x,x;y) &= y\jp(y\jp x^2) - x^2\jp y^2 + 2\,x\jp(y\jp A) - 2A^2,\\
 j(y,y,x;x) &= 2\,y\jp(x\jp A) - 2A^2 + x\jp(x\jp y^2) - x^2\jp y^2 .
\end{align*}
Hence
$0 = \tfrac12 j(y,x,x;y) - \tfrac12 j(y,y,x;x)
 = x\jp(y\jp A) - y\jp(x\jp A) - \tfrac12\bigl(x\jp(x\jp y^2) - y\jp(y\jp x^2)\bigr)$.
Since $\tfrac14(U_xy^2 - U_yx^2) = \tfrac14\bigl(2x\jp(x\jp y^2) - x^2\jp y^2 -
2y\jp(y\jp x^2) + y^2\jp x^2\bigr) = \tfrac12\bigl(x\jp(x\jp y^2) - y\jp(y\jp x^2)\bigr)$,
this is the claim.
\end{proof}

So the hypothesis of Theorem~\ref{thm:main} says exactly that $D(A)=0$.

\begin{lemma}[Fuglede--Putnam identity]\label{lem:3}
For all $x,y$: $\;2[L_x,L_{x\jp y}] + [L_y,L_{x^2}] = 0$. In particular, if
$x\opc x\jp y$ then $y\opc x^2$.
\end{lemma}
\prov{Their work: vdW 2020, Lemma 2.4(b)}{A standard linearisation; its reading as the Jordan shadow of the Fuglede--Putnam conclusion $x^2y = yx^2$ is the special-algebra step of \cite[Prop.~3.2]{vdW}.}
\begin{proof}
This is (J) with $(a,b,c) = (x,x,y)$: $[L_x,L_{x\jp y}] + [L_x,L_{y\jp x}] + [L_y,L_{x^2}]=0$.
\end{proof}

\begin{lemma}[intertwining]\label{lem:2}
For all $x,y$, with $A = x\jp y$ and $D=[L_x,L_y]$:
\[
 [L_A,U_x] + D\,U_x + U_x\,D = 0 .
\]
\end{lemma}
\prov{Literature: McCrimmon's Commuting Formula (FFII)}{With $V_{x,y} = 2(L_{x\jp y}+[L_x,L_y])$ this is $V_{x,y}U_x = U_xV_{y,x}$ \cite[5.2.3]{McC}, the quadratic Jordan axiom QJ2 \cite{McC66,Jac69}; only the hand derivation below from four instances of (J) is ours.}
\begin{proof}
Write $X = L_x$, $Y = L_y$, $\cP = L_A$, $\cQ = L_{x^2}$, $\cL = L_{x\jp A}$, so
$U_x = 2X^2-\cQ$ and $D = XY-YX$. Expanding,
\begin{align*}
 [\cP,U_x] &= 2\cP X^2 - 2X^2\cP + [\cQ,\cP],\\
 DU_x &= 2XYX^2 - XY\cQ - 2YX^3 + YX\cQ,\qquad
 U_xD = 2X^3Y - 2X^2YX - \cQ XY + \cQ YX,
\end{align*}
so the left side of the lemma, call it $\cK$, is
\begin{equation}\label{eq:Kexp}
 \cK = 2\cP X^2 - 2X^2\cP + [\cQ,\cP] + 2XYX^2 - 2YX^3 - XY\cQ + YX\cQ + 2X^3Y - 2X^2YX - \cQ XY + \cQ YX .
\end{equation}
Now consider the four operators
\begin{align*}
 J_1 &= [X,\cQ] &&= \tfrac13\,\text{(J) at } (x,x,x),\\
 J_2 &= 2[X,\cP] + [Y,\cQ] &&= \text{(J) at } (x,x,y),\\
 J_3 &= 2[X,\cL] + [\cP,\cQ] &&= \text{(J) at } (x,x,A),\\
 \cR &= \cL - 2\cP X - \cQ Y + X^2Y + YX^2 &&= \text{(II) at } (b,c,d) = (x,y,x),
\end{align*}
all of which vanish. (For $\cR$: (II) gives $L_{(x\jp y)\jp x} = L_{x\jp y}L_x +
L_{x\jp x}L_y + L_{y\jp x}L_x - L_xL_xL_y - L_yL_xL_x$.) We claim that
\begin{equation}\label{eq:Kid}
 \cK = -J_3 + 2[X,\cR] + J_1Y + YJ_1 - [X,J_2],
\end{equation}
as an identity of non-commutative polynomials in $X,Y,\cP,\cQ,\cL$; this proves the lemma.
To check~\eqref{eq:Kid}, expand the right side:
\begin{align*}
 -J_3 &= -2X\cL + 2\cL X + [\cQ,\cP],\\
 2[X,\cR] &= 2X\cL - 2\cL X - 4X\cP X + 4\cP X^2 - 2X\cQ Y + 2\cQ YX\\ &\qquad + 2X^3Y - 2X^2YX + 2XYX^2 - 2YX^3,\\
 J_1Y+YJ_1 &= X\cQ Y - \cQ XY + YX\cQ - Y\cQ X,\\
 -[X,J_2] &= -2X^2\cP + 4X\cP X - 2\cP X^2 - XY\cQ + X\cQ Y + Y\cQ X - \cQ YX .
\end{align*}
Adding: the $\cL$-terms cancel; $X\cP X$: $-4+4 = 0$; $\cP X^2$: $4-2=2$; $X\cQ Y$: $-2+1+1 = 0$;
$\cQ YX$: $2-1 = 1$; $Y\cQ X$: $-1+1 = 0$; the remaining terms
$[\cQ,\cP]$, $-2X^2\cP$, $2X^3Y$, $-2X^2YX$, $2XYX^2$, $-2YX^3$, $-\cQ XY$, $YX\cQ$, $-XY\cQ$ occur once
each. The sum is exactly~\eqref{eq:Kexp}.
\end{proof}

\begin{remark}
The derivation is how one finds~\eqref{eq:Kid}: replace $[\cQ,\cP]$ by $2[X,\cL]$ ($J_3$), then
$\cL$ by $2\cP X+\cQ Y-X^2Y-YX^2$ ($\cR$), then $\cQ X$ by $X\cQ$ ($J_1$). What is left is
$\cK = -2[X,[X,\cP]] + [X,[\cQ,Y]]$, and $J_2$ says $[\cQ,Y] = 2[X,\cP]$, so $\cK = 0$.
In element form~\eqref{eq:Kid} reads
$\cK (z) = -j(x,x,A;z) - 2j(x\jp z,x,y;x) + 2x\jp j(z,x,y;x) + \tfrac13 j(x,x,x;y\jp z)
+ \tfrac13 y\jp j(x,x,x;z) + j(x,x,y;x\jp z) - x\jp j(x,x,y;z)$. The formal proof uses a
different combination of ten instances,
$\cK (z) = \tfrac13 j(y,z,x;x^2) + \tfrac43 x\jp j(x,y,z;x) - \tfrac13 j(z,x,x;A)
- 2j(x\jp z,y,x;x) + \tfrac23 j(y,x,x;x\jp z) + \tfrac13 j(x\jp z,x,x;y)
- x\jp j(x,y,x;z) - \tfrac23 j(A,x,x;z) + \tfrac13 y\jp j(x,x,x;z) + \tfrac13 j(y\jp z,x,x;x)$.
Both combinations were checked to equal $\cK (z)$ exactly in the free commutative
non-associative algebra on $x,y,z$ (see the last section).
\end{remark}

\section*{Part II. The flow}

\begin{lemma}[exponentials of derivations]\label{lem:expder}
Let $\Delta$ be a bounded derivation of $\cA$ and $t\in\mathbb R$. Then $\Exp(t\Delta)$ is a
unital Jordan automorphism, it is positive, and $\|\Exp(t\Delta)v\|\le\|v\|$ for all $v$.
\end{lemma}
\prov{Literature (standard)}{Exponentials of bounded derivations are automorphisms in any Banach algebra; unital positive maps of order-unit spaces are contractive.}
\begin{proof}
(i) $\Delta(1) = \Delta(1\jp1) = 2\Delta(1)$, so $\Delta(1)=0$. The map
$s\mapsto\Exp(s\Delta)1$ has derivative $\Exp(s\Delta)\Delta(1) = 0$, so it is constant,
equal to $1$. (A differentiable map $\mathbb R\to$ Banach space with zero derivative is
constant, by the mean value inequality.)

(ii) Fix $u,v$, put $E_s = \Exp(s\Delta)$ and $G(s) = E_{-s}\bigl(E_su\jp E_sv\bigr)$.
The product is bounded bilinear and $\frac{d}{ds}E_sw = \Delta E_sw$, $\frac{d}{ds}E_{-s} =
-E_{-s}\Delta$, so
\[
 G'(s) = -E_{-s}\Delta\bigl(E_su\jp E_sv\bigr) + E_{-s}\bigl(\Delta E_su\jp E_sv
 + E_su\jp\Delta E_sv\bigr) = 0
\]
because $\Delta$ is a derivation. Hence $G(t) = G(0) = u\jp v$; applying $E_t$ and using
$E_tE_{-t} = \Exp(t\Delta - t\Delta) = I$ gives $E_t(u\jp v) = E_tu\jp E_tv$.

(iii) If $v\ge0$ then $v = (\sqrt v)^2$, so $E_tv = (E_t\sqrt v)^2\ge 0$; thus $E_t$ is
positive, hence monotone.

(iv) Apply the monotone unital map $E_t$ to $-\|v\|1\le v\le\|v\|1$: we get
$-\|v\|1\le E_tv\le\|v\|1$, so $\|E_tv\|\le\|v\|$ by Fact~\ref{f:norm}.
\end{proof}

\begin{proposition}[the flow]\label{prop:flow}
Let $x,y\in\cA$ with $D(A) = 0$, where $A = x\jp y$, $D = [L_x,L_y]$. Then for every
$t\in\mathbb R$
\begin{equation}\label{eq:flow}
 \Exp(2tD)\Bigl(U_{e^{tA}}\bigl(U_x(e^{-2tA})\bigr)\Bigr) = x^2 ,
\end{equation}
and consequently $W(t) := U_{e^{tA}}U_x(e^{-2tA})$ satisfies $\|W(t)\|\le\|x^2\|$ for all
$t\in\mathbb R$.
\end{proposition}
\prov{New?}{A Jordan analogue of Rosenblum's exponential proof of Fuglede--Putnam \cite{Ros58} (in $B(H)$ it is $e^{\sigma A}xe^{-\sigma A}$ against $e^{-\sigma C/2}x^2e^{\sigma C/2}$, $C = [x,y]$); no Jordan Fuglede--Putnam theorem found in the literature.}
\begin{proof}
Put $U = U_x$ and $\Lambda = L_A$ (in this proof).

\emph{$\Lambda$ and $D$ commute.} By Lemma~\ref{lem:der}, $[D,\Lambda] = L_{D(A)} = 0$.

\emph{An intertwining.} Put $\Omega_+ = 2D+2\Lambda$ and $\Omega_- = 2D-2\Lambda$. Then
$\Omega_+U + U\Omega_- = 2\bigl(\Lambda U - U\Lambda + DU + UD\bigr) = 0$ by Lemma~\ref{lem:2}.

\emph{A constant function.} Let $F(s) = \Exp(s\Omega_+)\,U\,\Exp(s\Omega_-)\in B(\cA)$. By the product
rule in $B(\cA)$, with $\frac{d}{ds}\Exp(s\Omega_+) = \Exp(s\Omega_+)\Omega_+$ and $\frac{d}{ds}\Exp(s\Omega_-) =
\Omega_-\Exp(s\Omega_-)$,
\begin{align*}
 F'(s) &= \Exp(s\Omega_+)\Omega_+U\Exp(s\Omega_-) + \Exp(s\Omega_+)U\Omega_-\Exp(s\Omega_-)\\
 &= \Exp(s\Omega_+)(\Omega_+U+U\Omega_-)\Exp(s\Omega_-) = 0 .
\end{align*}
So $F(t) = F(0) = U$ for every $t$.

\emph{Splitting the exponentials.} Since $2tD$ and $2t\Lambda$ commute,
$\Exp(t\Omega_+) = \Exp(2tD)\Exp(2t\Lambda)$ and $\Exp(t\Omega_-) = \Exp(-2t\Lambda)\Exp(2tD)$.

\emph{Evaluation at $1$.} Apply $F(t) = U$ to $1$. First $\Exp(2tD)1 = 1$ by
Lemma~\ref{lem:expder}(i), since $2D$ is a derivation (Lemma~\ref{lem:der}). Next
$\Exp(-2t\Lambda)1 = \Exp(L_{-2tA})1 = e^{-2tA}$ (Fact~\ref{f:cfc}), and
$\Exp(2t\Lambda) = \Exp(2L_{tA}) = U_{e^{tA}}$ (Fact~\ref{f:ff}). Hence
\[
 \Exp(2tD)\,U_{e^{tA}}\,U_x\bigl(e^{-2tA}\bigr) = F(t)(1) = U_x(1) = x^2 .
\]
\emph{The bound.} Since $\Exp(-2tD)\Exp(2tD) = I$, \eqref{eq:flow} gives
$W(t) = \Exp(-2tD)(x^2)$, and $\|W(t)\|\le\|x^2\|$ by Lemma~\ref{lem:expder}(iv) applied to
the derivation $2D$ at time $-t$.
\end{proof}

\section*{Part III. The gap estimate and spectral projections}

For $x,h\in\cA$ and $\tau\in\mathbb R$ put $W_{x,h}(\tau) = U_{e^{\tau h}}U_x(e^{-2\tau h})$.
Proposition~\ref{prop:flow} says $W_{x,A}$ is bounded on all of $\mathbb R$ when $D(A)=0$.
The next lemma uses only continuous functions of $h$.

\begin{lemma}[gap lemma]\label{lem:gap}
Let $x,h\in\cA$, $C\ge0$, and suppose $\|W_{x,h}(\tau)\|\le C$ for all $\tau>0$. Let
$\alpha<\beta$ and $k,f\in C(\spec h)$ with $f\ge0$, $k(t) = 0$ for $t<\beta$ and $f(t)=0$
for $t>\alpha$. Then $Z := U_{k(h)}U_x f(h) = 0$.
\end{lemma}
\prov{New? (in this form)}{That a bounded exponential orbit kills off-diagonal spectral blocks is a standard operator-theoretic idea; the Jordan version with $U$-compressions and continuous cut-offs we have not found.}
\begin{proof}
First, $Z\ge 0$: $f(h)\ge0$ and $U_x$, $U_{k(h)}$ are positive (Facts~\ref{f:cfc},
\ref{f:pos}). Fix $\tau>0$.

\emph{Step 1: $f(h)\le\|f\|e^{2\tau\alpha}e^{-2\tau h}$.} By Fact~\ref{f:cfc} it suffices to
check $f(t)\le\|f\|e^{2\tau\alpha}e^{-2\tau t}$ on $\spec h$. If $t>\alpha$ the left side
is $0$ and the right side is $\ge 0$. If $t\le\alpha$ then $2\tau(\alpha - t)\ge0$, so
$e^{2\tau\alpha}e^{-2\tau t}\ge1$ and $f(t)\le\|f\|\le\|f\|e^{2\tau\alpha}e^{-2\tau t}$.

\emph{Step 2: positivity.} Since $U_x$ and $U_{k(h)}$ are monotone,
$0\le Z\le \|f\|e^{2\tau\alpha}\,U_{k(h)}U_x(e^{-2\tau h})$.

\emph{Step 3: factorisation.} Let $g(t) = e^{\tau t/2}$ and $m(t) = k(t)e^{-\tau t}$, so
$k = g^2m$ and $g^2(h) = e^{\tau h}$. Put $\gamma = g(h)$, $\mu = m(h)$. By Fact~\ref{f:Uopc},
$k(h) = \gamma^2\jp \mu = U_\gamma\mu$, so by the fundamental formula
$U_{k(h)} = U_\gamma U_\mu U_\gamma$. The elements $\gamma, \mu, \gamma^2, \mu^2$ lie in $C(h)$ and pairwise
operator-commute, so $L_\gamma, L_{\gamma^2}$ commute with $L_\mu, L_{\mu^2}$ and hence $U_\gamma$ commutes
with $U_\mu$. With $U_\gamma U_\gamma = U_{\gamma^2}$ (Fact~\ref{f:ff}) we get
$U_{k(h)} = U_\mu U_{\gamma^2} = U_{m(h)}U_{e^{\tau h}}$. Hence
$U_{k(h)}U_x(e^{-2\tau h}) = U_{m(h)}W_{x,h}(\tau)$.

\emph{Step 4: the norm of $m(h)$.} $\|m(h)\| = \sup_{t\in\spec h}|k(t)|e^{-\tau t}$. For
$t<\beta$ the term is $0$; for $t\ge\beta$ it is $\le\|k\|e^{-\tau\beta}$. So
$\|m(h)\|\le\|k\|e^{-\tau\beta}$.

\emph{Step 5: conclusion.} By Steps 2--4 and Fact~\ref{f:norm},
\[
 \|Z\|\le\|f\|e^{2\tau\alpha}\,\|U_{m(h)}W_{x,h}(\tau)\|
 \le\|f\|e^{2\tau\alpha}\cdot 3\|k\|^2e^{-2\tau\beta}\cdot C
 = 3C\|f\|\|k\|^2\,e^{-2\tau(\beta-\alpha)} .
\]
The left side does not depend on $\tau$, and the right side tends to $0$ as
$\tau\to\infty$ because $\beta-\alpha>0$. So $\|Z\| = 0$.
\end{proof}

Note that no spectral projection has to commute with $U_{e^{\tau h}}$: the only
operators moved past each other are $U$'s of elements of $C(h)$.

\begin{proposition}[passage to spectral projections, JBW]\label{prop:stepB}
Let $\cA$ be a JBW-algebra and $x,h\in\cA$ with $\|W_{x,h}(\tau)\|\le C$ for all $\tau>0$.
Then $\Pp(p)x = 0$, hence $p\opc x$, for every spectral projection $p = p_j$ of $h$ as in
Fact~\ref{f:jbw} (every $\delta>0$, $j\in\mathbb N$).
\end{proposition}
\prov{New? (in this form)}{Combines the gap lemma with normality \cite[\S4.1]{HOS} and Facts~\ref{f:S3}, \ref{f:peirce}; not found.}
\begin{proof}
Fix $\delta,j$, let $\rho = \rho_j$ and $s = j\delta - \|h\|$, so $\rho(t) = (t-s)^+$, and
$p = \pi_\rho = \sup_n\sigma_n(h)$. By Lemma~\ref{lem:peirce} it suffices to
show $U_pU_x(1-p) = 0$. Put $u = U_x(1-p)$; $u\ge0$ since $0\le p\le 1$. Let $r = \sqrt u$.

\emph{(a) $U_r k_m(h) = 0$ for $k_m = (\rho - \tfrac{2}{m+1})^+$, $m\in\mathbb N$.} Let
$f_m = 1-\sigma_m = 1 - \min((m+1)\rho, 1)\ge 0$. Then $f_m(t) = 0$ for
$t> \alpha := s+\tfrac1{m+1}$ (there $(m+1)\rho(t)>1$), and $\sqrt{k_m}(t) = 0$ for
$t<\beta := s+\tfrac2{m+1}$ (there $\rho(t)<\tfrac2{m+1}$). Lemma~\ref{lem:gap} with
$k = \sqrt{k_m}$, $f = f_m$ gives $U_{\sqrt{k_m}(h)}U_xf_m(h) = 0$. Since
$\sigma_m(h)\le p$ we have $1-p\le f_m(h)$, so by monotonicity
$0\le U_{\sqrt{k_m}(h)}u\le U_{\sqrt{k_m}(h)}U_xf_m(h) = 0$. As
$\sqrt{k_m(h)} = \sqrt{k_m}(h)$, Fact~\ref{f:S3} (with $a = k_m(h)$, $b = u$) gives
$U_{r}k_m(h) = 0$.

\emph{(b) $U_r\rho(h) = 0$.} $\|k_m-\rho\|_\infty\le\tfrac2{m+1}$, since
$\rho - \tfrac2{m+1}\le k_m\le\rho$. By Fact~\ref{f:cfc}, $k_m(h)\to\rho(h)$ in norm, and
$U_r$ is continuous, so $U_r\rho(h) = \lim U_rk_m(h) = 0$.

\emph{(c) $U_rp = 0$.} Since $\sigma_n\le(n+1)\rho$, monotonicity gives
$0\le U_r\sigma_n(h)\le(n+1)U_r\rho(h) = 0$. The $\sigma_n(h)$ increase to $p$ inside
$[0,1]$, so by normality $U_rp = \sup_nU_r\sigma_n(h) = 0$.

\emph{(d) Conclusion.} Fact~\ref{f:S3} with $(a,b) = (u,p)$ gives $U_{\sqrt p}u = 0$. As
$p$ is an idempotent, $\sqrt p = p$, so $U_pU_x(1-p) = U_pu = 0$. By
Lemma~\ref{lem:peirce}, $\Pp(p)x = 0$, and by Fact~\ref{f:peirce}, $p\opc x$.
\end{proof}

\begin{proposition}[from spectral projections to powers]\label{prop:pows}
Let $\cA$ be a JBW-algebra, $h,z\in\cA$, and suppose $p_j\opc z$ for all spectral projections
$p_j$ of $h$ (all $\delta>0$, $j$). Then $h\opc z$, and $h^n\opc z$ for every $n$;
hence $z\in\comm(h)$.
\end{proposition}
\prov{Literature (standard)}{Spectral approximation by idempotents, cf.\ \cite[4.2.3]{HOS} and the same use in \cite[Props.~2.33--2.35]{vdW}.}
\begin{proof}
Fix $\delta>0$. Let $S_\delta$ be the linear span of $1$ and the $p_j$. For $i\le j$ we
have $p_j\le p_i$, so $p_i\jp p_j = p_j = p_{\max(i,j)}$ (Fact~\ref{f:peirce}); thus
$S_\delta$ is closed under the product, and every power of an element of $S_\delta$ lies
in $S_\delta$. Every element of $S_\delta$ operator-commutes with $z$ (linearity).
Let $h_\delta = c1+\delta\sum_{i<m}p_{i+1}\in S_\delta$ be the approximant
of~\eqref{eq:specapprox}. Then $h_\delta^n\opc z$ and $\|h-h_\delta\|\le\delta$. Letting
$\delta = 1/(k+1)\to0$, $h_\delta^n\to h^n$ by continuity of $w\mapsto w^n$, and the set
$\{c: c\opc z\}$ is closed, so $h^n\opc z$ (the case $n=1$ gives $h\opc z$).
Consequently $q(h)\opc z$ for every polynomial $q$, and since polynomials in $h$ are dense
in $C(h)$, every $c\in C(h)$ satisfies $c\opc z$.
\end{proof}

\section*{Part IV. The Fuglede--Putnam step}

\begin{proposition}\label{prop:ya}
Let $\cA$ be a JBW-algebra, $a,b\ge0$ with $U_{\sqrt a}b = U_{\sqrt b}a$. Then $y\opc a$.
\end{proposition}
\prov{New? (proof); statement from their work}{Follows from \cite[Thm.~3.12]{vdW}, which uses Shirshov--Cohn and Shultz's structure theorem; an intrinsic (structure-free) proof we have not found.}
\begin{proof}
By Lemma~\ref{lem:1}, $D(A) = \tfrac14(U_xy^2-U_yx^2) = \tfrac14(U_{\sqrt a}b -
U_{\sqrt b}a) = 0$. By Proposition~\ref{prop:flow}, $\|W_{x,A}(\tau)\|\le\|x^2\|$ for all
$\tau$. By Proposition~\ref{prop:stepB} (with $h = A$), $x$ operator-commutes with every
spectral projection of $A$, so $A\opc x$ by Proposition~\ref{prop:pows}. By
Lemma~\ref{lem:3}, $y\opc x^2 = a$.
\end{proof}

In a special algebra the last step is the familiar
$x(xy+yx) = (xy+yx)x\Rightarrow x^2y = yx^2$.

\section*{Part V. Conclusion in JBW-algebras}

\begin{proposition}[the $(y+\lambda,a)$ trick]\label{prop:trick}
Let $\cA$ be a JBW-algebra and $y\opc a$. Then $a^n\opc y$ for every $n$, so $y\in\comm(a)$.
\end{proposition}
\prov{New? (minor)}{A device of this note (and its formalisation); the conclusion also follows from \cite[Thm.~3.12]{vdW}.}
\begin{proof}
Fix $\lambda = m+1$, $m\in\mathbb N$, and put $x_\lambda = y+\lambda1$. Then
$x_\lambda\opc a$ (as $y\opc a$ and $1\opc a$), so $D_\lambda := [L_{x_\lambda},L_a] = 0$
and in particular $D_\lambda(x_\lambda\jp a) = 0$. Proposition~\ref{prop:flow} for the pair
$(x_\lambda,a)$ shows that $A_\lambda = x_\lambda\jp a$ has
$\|W_{x_\lambda,A_\lambda}(t)\|\le\|x_\lambda^2\|$ for all $t$. Let
$h_\lambda = \lambda^{-1}A_\lambda = a + \lambda^{-1}y\jp a$. Since $\tau h_\lambda =
(\tau/\lambda)A_\lambda$ and $-2\tau h_\lambda = -2(\tau/\lambda)A_\lambda$,
$W_{x_\lambda,h_\lambda}(\tau) = W_{x_\lambda,A_\lambda}(\tau/\lambda)$, which is bounded by
$\|x_\lambda^2\|$. By Proposition~\ref{prop:stepB}, $x_\lambda$ operator-commutes with every
spectral projection of $h_\lambda$, hence so does $y = x_\lambda-\lambda1$. By
Proposition~\ref{prop:pows}, $h_\lambda^n\opc y$ for all $n$. Now
$\|h_\lambda - a\| = \lambda^{-1}\|y\jp a\|\to0$ as $m\to\infty$, so $h_\lambda^n\to a^n$,
and by closedness $a^n\opc y$. The last clause is Proposition~\ref{prop:pows}'s final step.
\end{proof}

\begin{proposition}[from $y$ to $b=y^2$, JBW]\label{prop:mulmem}
In a JBW-algebra, if $y\in\comm(c)$ then $y^2\in\comm(c)$.
\end{proposition}
\prov{Their work: vdW 2020, Prop.~2.34}{Same route: idempotents and Jacobson's subalgebra lemma \cite{Jac89}.}
\begin{proof}
Let $w\in C(c)$ and let $q$ be a spectral projection of $w$. The elements $\sigma_n(w)$
lie in $C(w)\subseteq C(c)$, so they operator-commute with $y$, and by normality so does
their supremum $q$. By Fact~\ref{f:peirce}, $\Pp(q)y = 0$, i.e.\ $y\in\cA_1(q)+\cA_0(q)$,
which is a subalgebra; so $\Pp(q)(y^2) = 0$ and $q\opc y^2$. This holds for every spectral
projection of $w$, so $w\opc y^2$ by Proposition~\ref{prop:pows}.
\end{proof}

\begin{proof}[Proof of Theorem~\ref{thm:main} for JBW-algebras]
By Proposition~\ref{prop:ya}, $y\opc a$; by Proposition~\ref{prop:trick}, $y\in\comm(a)$;
by Proposition~\ref{prop:mulmem}, $b = y^2\in\comm(a)$.
\end{proof}

The next lemma is not needed for JBW-algebras, but it is the tool that replaces
Proposition~\ref{prop:mulmem} and the $(y+\lambda,a)$ trick in Part VI.

\begin{lemma}[(SQ): $u\opc v\Rightarrow v\opc u^2$]\label{lem:SQ}
Suppose that in $\cA$ every pair $x,h$ with $W_{x,h}$ bounded on $\mathbb R$ satisfies
$x\opc h$ (true in JBW-algebras by Propositions~\ref{prop:stepB} and~\ref{prop:pows}, and
in all JB-algebras by Part VI). Then $u\opc v$ implies $v\opc u^2$.
\end{lemma}
\prov{Statement: their work; proof: New?}{The implication is part of \cite[Thm.~3.13, (a)$\Leftrightarrow$(d)]{vdW}, obtained there from structure theory (cf.\ his Remark~2.15); this flow-based proof we have not found.}
\begin{proof}
Let $D' = [L_u,L_v]$; it is $0$ by hypothesis, so $D'(u\jp v) = 0$. Proposition~\ref{prop:flow}
for the pair $(u,v)$ gives $\|W_{u,u\jp v}(t)\|\le\|u^2\|$ for all $t\in\mathbb R$. By
assumption $u\opc u\jp v$, and Lemma~\ref{lem:3} (with $(x,y) = (u,v)$) gives $v\opc u^2$.
\end{proof}

\section*{Part VI. JB-algebras}

In a JB-algebra there need be no spectral projections. We isolate what Parts III--V used
them for.

\begin{proposition}[continuous Peirce criterion]\label{prop:cp}
Let $\cA$ be a JB-algebra and $x,h\in\cA$. Suppose (G): $U_{k(h)}U_xf(h) = 0$ for all
$k,f\in C(\spec h)$ with $f\ge0$ and \emph{separated supports}, meaning that for some
$\alpha<\beta$ either [$k=0$ on $\{t<\beta\}$ and $f=0$ on $\{t>\alpha\}$] or
[$k=0$ on $\{t>\alpha\}$ and $f=0$ on $\{t<\beta\}$]. Then $x\opc h$.
\end{proposition}
\prov{New?}{A spectral-projection-free replacement for Props.~\ref{prop:stepB}--\ref{prop:pows}; not found. Its proof (VI.1--VI.7) is annotated step by step.}

\begin{lemma}\label{lem:twosided}
If $\|W_{x,h}(\tau)\|\le C$ for all $\tau\in\mathbb R$, then (G) holds.
\end{lemma}
\prov{New? (as Lemma~\ref{lem:gap})}{Two applications of the gap lemma.}
\begin{proof}
The first order is Lemma~\ref{lem:gap}. For the second, apply Lemma~\ref{lem:gap} to
$-h$ in place of $h$: $W_{x,-h}(\tau) = W_{x,h}(-\tau)$ is bounded for $\tau>0$;
$\spec(-h) = -\spec h$, and $\tilde k(t) = k(-t)$, $\tilde f(t) = f(-t)$ satisfy
$\tilde k(-h) = k(h)$, $\tilde f(-h) = f(h)$, $\tilde f\ge0$, $\tilde k = 0$ on
$\{t<-\alpha\}$ and $\tilde f = 0$ on $\{t>-\beta\}$, with $-\beta<-\alpha$. (The formal
proof instead carries a sign $s=\pm1$ through the proof of Lemma~\ref{lem:gap}, replacing
$\tau$ by $s\tau$ and $t$ by $st$; the estimates are identical.)
\end{proof}

The remainder of this part proves Proposition~\ref{prop:cp}. Fix $x,h$ with (G).

\paragraph{VI.1 Multiplicativity of $U$ on $C(h)$.}
For $s,s_1,s_2\ge 0$ and $m,m_1,m_2$ in $C(\spec h)$ (all evaluated at $h$):
\begin{align*}
 \text{(M0)}\quad & U_{sm} = U_mU_s,\\
 \text{(Ma)}\quad & U_{sm_1,sm_2} = U_{m_1,m_2}U_s,\\
 \text{(Mb)}\quad & U_{s_1m_1,s_2m_2} + U_{s_2m_1,s_1m_2} = 2\,U_{m_1,m_2}U_{s_1,s_2}.
\end{align*}
(M0) is Step 3 of Lemma~\ref{lem:gap} with $g = \sqrt s$. (Ma): apply (M0) to $m_1+m_2$,
expand both sides with $U_{a+c} = U_a+U_c+2U_{a,c}$, and cancel (M0) for $m_1$ and
$m_2$. (Mb): apply (Ma) with $s = s_1+s_2\ge0$, expand in $s$ in the same way, and cancel
(Ma) for $s_1$ and for $s_2$. Two consequences, for $p\cdot r = 0$ pointwise:
\prov{Literature (routine)}{Multiplicativity of $U$ on an associative subalgebra from the fundamental formula; the polarisations (Ma), (Mb) are routine.}
\begin{itemize}
\item if $r\ge0$ then $U_{p(h)}U_{r(h)} = 0$: (M0) with $s = r$, $m = p$ gives
 $U_{rp} = U_pU_r$, and $rp = 0$;
\item if $r\ge0$ then $U_{p(h)}L_{r(h)} = 0$: (Mb) with $s_1 = 1$, $s_2 = r$,
 $m_1=m_2=p$ gives $U_{p,rp}+U_{rp,p} = 2U_pU_{1,r} = 2U_pL_r$, and $rp=0$.
\end{itemize}

\paragraph{VI.2 Orthogonality lemma.} \emph{If $e\ge0$ and $U_ce = 0$ then $c\jp e = 0$.}
For every $t\in\mathbb R$,
$0\le U_{1+tc}e = e + 2t\,c\jp e + t^2U_ce = e+2t\,c\jp e$ (using
$U_{1+tc} = I + 2tL_c + t^2U_c$, $I$ the identity operator). For $s>0$, taking $t = \mp s$ and dividing by $2s$ gives
$-\tfrac1{2s}e\le c\jp e\le\tfrac1{2s}e$, so
$-\tfrac{\|e\|}{2s}1\le c\jp e\le\tfrac{\|e\|}{2s}1$ and $\|c\jp e\|\le\|e\|/(2s)$ for every
$s>0$. Hence $c\jp e = 0$.
\prov{Literature (variant)}{For $c\ge 0$ this is in Alfsen--Shultz \cite[Lemma~1.26]{AS}; the $U_{1+tc}$ argument is routine.}

\paragraph{VI.3 A cross term of the fundamental formula.} \emph{If $a\jp b = 0$ then for all $x$}
\[
 4\,(U_{a,b}x)^2 + 2\,(U_ax)\jp(U_bx) = U_aU_x(b^2) + U_bU_x(a^2).
\]
Applying the fundamental formula to $1$ gives $(U_cx)^2 = U_cU_x(c^2)$ for every $c$. Use it
for $c = a+b$ and $c = a-b$; since $a\jp b = 0$, $(a\pm b)^2 = a^2+b^2$, and
$U_{a\pm b} = U_a+U_b\pm2U_{a,b}$. Averaging the two equations,
$(U_ax+U_bx)^2 + 4(U_{a,b}x)^2 = (U_a+U_b)U_x(a^2+b^2)$. Subtracting
$(U_ax)^2 = U_aU_xa^2$ and $(U_bx)^2 = U_bU_xb^2$ gives the claim.
\prov{Literature (routine)}{A polarisation of FF applied to $1$.}

\paragraph{VI.4 Off-diagonal parts vanish.} \emph{If $p,r\ge0$ in $C(\spec h)$ have
separated supports, then $U_{\hat p,\hat r}\,x = 0$, where $\hat p = p(h)$, $\hat r = r(h)$.}
We have $pr = 0$, so $\hat p\jp \hat r = 0$. The pairs $(p,r^2)$ and $(r,p^2)$ also have separated
supports (the second in the opposite order), so (G) gives
$U_{\hat p}U_x(\hat r^2) = 0$ and $U_{\hat r}U_x(\hat p^2) = 0$. Put $l = \|x\|$ and $x' = x+l1\ge 0$. Since
$U_{x'}v = U_xv + 2l\,x\jp v + l^2v$,
\[
 U_{\hat p}U_{x'}(\hat r^2) = U_{\hat p}U_x(\hat r^2) + 2l\,U_{\hat p}L_{\hat r^2}x + l^2\,U_{\hat p}U_{\hat r}(1) = 0
\]
by VI.1: the second bullet (with $r^2\ge0$ in place of $r$, and $pr^2 = 0$) gives
$U_{\hat p}L_{\hat r^2} = 0$, and the first bullet (with $r\ge0$, $pr = 0$) gives
$U_{\hat p}U_{\hat r} = 0$. Likewise $U_{\hat r}U_{x'}(\hat p^2) = 0$. Also
$U_{\hat p,\hat r}1 = \hat p\jp \hat r = 0$, so $z := U_{\hat p,\hat r}x' = U_{\hat p,\hat r}x$. Let $c = U_{\hat p}x'$ and $e = U_{\hat r}x'\ge0$.
By the fundamental formula and the first bullet of VI.1, $U_ce = U_{\hat p}U_{x'}U_{\hat p}U_{\hat r}(x') = 0$, so $c\jp e = 0$ by
VI.2. Now VI.3 (with $a = \hat p$, $b = \hat r$, and $x'$) gives $4z^2 = 0$, so
$\|z\|^2 = \|z^2\| = 0$.
\prov{New?}{Off-diagonal vanishing from (G) without spectral projections; not found.}

\paragraph{VI.5 Second-order vanishing: $h^2\jp x = h\jp(h\jp x)$.}
Let $\Delta = L_{h^2}-L_h^2$. For $\mu\in\mathbb R$ and $g = h-\mu1$ one checks
$\Delta = L_{g^2}-L_g^2 = \tfrac12\bigl(U_{1,g^2} - U_g\bigr)$ (expand $g$; and
$U_{1,g^2} = L_{g^2}$, $U_g = 2L_g^2 - L_{g^2}$).

\emph{Ramps.} For $\varepsilon>0$ let $\mathrm r_c(t) = \min\bigl(1,\max(0,(t-c)/\varepsilon)\bigr)$;
it is $0$ for $t\le c$, $1$ for $t\ge c+\varepsilon$, and decreasing in $c$.

\emph{One piece: $\|\Delta(\phi(h)\jp x)\|\le12\varepsilon^2\|x\|$ for
$\phi = \mathrm r_c-\mathrm r_{c+\varepsilon}$.} Let $\psi = \mathrm r_{c-2\varepsilon}
-\mathrm r_{c+3\varepsilon}$, $\ell o = 1-\mathrm r_{c-2\varepsilon}$,
$hi = \mathrm r_{c+3\varepsilon}$, so $\psi+\ell o+hi = 1$ and all four functions are
$\ge0$, with $0\le\phi,\psi\le1$. $\phi$ vanishes outside $[c,c+2\varepsilon]$, $\psi$ outside
$[c-2\varepsilon,c+4\varepsilon]$, $\ell o$ on $\{t\ge c-\varepsilon\}$, and $hi$ on
$\{t\le c+3\varepsilon\}$. So $(\phi,\ell o)$ are separated ($\alpha = c-\varepsilon$,
$\beta = c$) and $(\phi,hi)$ are separated ($\alpha = c+2\varepsilon$, $\beta =
c+3\varepsilon$). With $\Phi = \phi(h)$, $\Psi = \psi(h)$ and VI.4,
\[
 \Phi\jp x = U_{\Phi,1}x = U_{\Phi,\Psi}x + U_{\Phi,\ell o(h)}x + U_{\Phi,hi(h)}x = U_{\Phi,\Psi}x.
\]
Take $\mu = c+\varepsilon$, $g = h-\mu1$. (Mb) with $(s_1,s_2) = (\phi,\psi)$ and
$(m_1,m_2) = (1,g^2)$, resp.\ $(g,g)$, gives
$U_{\phi,\psi g^2}+U_{\psi,\phi g^2} = 2U_{1,g^2}U_{\phi,\psi}$ and
$U_{\phi g,\psi g}+U_{\psi g,\phi g} = 2U_gU_{\phi,\psi}$. Therefore
\[
 \Delta(\Phi\jp x) = \tfrac12(U_{1,g^2}-U_g)U_{\Phi,\Psi}x
 = \tfrac14\bigl(U_{\phi,\psi g^2}+U_{\psi,\phi g^2}\bigr)x
 - \tfrac14\bigl(U_{\phi g,\psi g}+U_{\psi g,\phi g}\bigr)x .
\]
On the support of $\phi$, $|t-\mu|\le\varepsilon$; on that of $\psi$, $|t-\mu|\le
3\varepsilon$. So $\|\phi\|\le1$, $\|\psi\|\le 1$, $\|\phi g\|\le\varepsilon$,
$\|\psi g\|\le3\varepsilon$, $\|\phi g^2\|\le\varepsilon^2$, $\|\psi g^2\|\le 9\varepsilon^2$
(sup norms over $\spec h$, equal to the norms at $h$). By
$\|U_{a,b}w\|\le3\|a\|\|b\|\|w\|$,
\[
 \|\Delta(\Phi\jp x)\|\le\tfrac14\bigl(3\cdot9\varepsilon^2 + 3\varepsilon^2\bigr)\|x\|
 + \tfrac14\bigl(9\varepsilon^2+9\varepsilon^2\bigr)\|x\| = 12\,\varepsilon^2\|x\| .
\]

\emph{Summing.} Let $R_0 = \|h\|$, $m_0 = \lceil 2R_0+2\rceil$, $n\ge1$, $\varepsilon = 1/n$,
$c_i = -R_0-1+i\varepsilon$, $f_i = \mathrm r_{c_i}$ and $N = m_0n$. On
$\spec h\subseteq[-R_0,R_0]$: $f_0 = 1$ because $t\ge -R_0\ge c_0+\varepsilon$
($\varepsilon\le1$), and $f_N = 0$ because $c_N = -R_0-1+m_0\ge R_0+1> t$. So the
telescoping sum gives $\sum_{i<N}(f_i-f_{i+1}) = f_0 - f_N = 1$ on $\spec h$, and
$x = \sum_{i<N}(f_i-f_{i+1})(h)\jp x$, each $f_i-f_{i+1} = \mathrm r_{c_i}-\mathrm
r_{c_i+\varepsilon}$ being a piece. By the triangle inequality,
\[
 \|\Delta x\|\le N\cdot12\varepsilon^2\|x\| = \frac{12m_0\|x\|}{n}\xrightarrow[n\to\infty]{}0 .
\]
Hence $\Delta x = h^2\jp x - h\jp(h\jp x) = 0$. No estimate for sums of orthogonal pieces
is needed: $O(1/\varepsilon)$ pieces of size $O(\varepsilon^2)$ suffice.
\prov{New?}{The ramp decomposition and $12\varepsilon^2$ estimate; not found.}

\paragraph{VI.6 Kleinecke--Shirokov.} Let $\mathfrak D = [L_x,L_h]$, a derivation
(Lemma~\ref{lem:der}). Then $\mathfrak D(h) = x\jp h^2 - h\jp(x\jp h) = \Delta x = 0$, so
$[\mathfrak D,L_h] = L_{\mathfrak D(h)} = 0$. In the Banach algebra $B(\cA)$ put $T = L_x$, $S = L_h$ and
$\delta(X) = XS-SX$. Then $\delta(T) = \mathfrak D$, $\delta$ is a derivation of $B(\cA)$, and
$\delta(\mathfrak D) = \mathfrak DS-S\mathfrak D = 0$. \emph{Claim: $n!\,\|\mathfrak D^n\|\le(2\|S\|\|T\|)^n$.}
First, by induction on $m$, for every $Y$,
\[
 \delta^{m+1}(TY) = T\,\delta^{m+1}(Y) + (m+1)\,\mathfrak D\,\delta^m(Y):
\]
for $m=0$ this is $\delta(TY) = \mathfrak D Y+T\delta(Y)$; the step applies $\delta$, using
$\delta(\mathfrak D\,\delta^m Y) = \mathfrak D\,\delta^{m+1}Y$. Next, by induction on $n$,
$\delta^n(T^n) = n!\,\mathfrak D^n$ and $\delta^{n+1+k}(T^n) = 0$ for all $k$: for $n=0$,
$\delta(1) = 0$; for the step, $T^{n+1} = T\cdot T^n$ and the display give
$\delta^{n+1}(T^{n+1}) = T\delta^{n+1}(T^n) + (n+1)\mathfrak D\,\delta^n(T^n) = (n+1)!\,\mathfrak D^{n+1}$ and
$\delta^{n+2+k}(T^{n+1}) = T\delta^{n+2+k}(T^n)+(n+2+k)\mathfrak D\,\delta^{n+1+k}(T^n) = 0$.
Finally $\|\delta(X)\|\le2\|S\|\|X\|$, so
$n!\|\mathfrak D^n\| = \|\delta^n(T^n)\|\le(2\|S\|)^n\|T\|^n$. With $K = 2\|L_h\|\|L_x\|$:
$n!\,\|\mathfrak D^nv\|\le K^n\|v\|$ for all $n$, $v$.
\prov{Literature: Kleinecke 1957, Shirokov 1956}{The bound $\|\delta^n(T^n)\| = n!\|\mathfrak D^n\|$ is the standard Kleinecke--Shirokov argument \cite{Kle57,Shi56}; applied here in $B(\cA)$.}

\paragraph{VI.7 Gelfand--Hille.} \emph{Let $E$ be a real Banach space and $\mathfrak D\in B(E)$ with
$n!\|\mathfrak D^nv\|\le K^n\|v\|$ and $\|\Exp(t\mathfrak D)v\|\le\|v\|$ for all $n$, $v$, and real $t$. Then
$\mathfrak D = 0$.}
Fix $v$ and a functional $\ell$; put $B = \|\ell\|\|v\|$ and $c_n = \ell(\mathfrak D^nv)$, so
$|c_n|\le BK^n/n!$. For $r\ge0$,
\[
 \frac{|c_n|}{n!}r^n\le B\frac{(Kr)^n}{(n!)^2} = B\Bigl(\frac{\sqrt{Kr}^{\,n}}{n!}\Bigr)^2 .
\]
For $s\ge0$, $s^n/n!\le e^s$, so $(s^n/n!)^2\le e^s s^n/n!$ and
$\sum_n(s^n/n!)^2\le e^{2s}$. Hence $F(w) = \sum_nc_nw^n/n!$ converges absolutely and
locally uniformly on $\mathbb C$, $F$ is entire, and
$|F(w)|\le B e^{2\sqrt{K|w|}}$. For real $s$, $F(s) = \ell(\Exp(s\mathfrak D)v)$ (apply $\ell$ to the
exponential series), so $|F(s)|\le\|\ell\|\|\Exp(s\mathfrak D)v\|\le B$.
Let $\Gamma(z) = F(z^2)$, entire with $|\Gamma(z)|\le Be^{2\sqrt K|z|}$. On the real axis $z^2\ge0$
and on the imaginary axis $z^2 = -|z|^2\le 0$ are real, so $|\Gamma|\le B$ on both axes. In each
closed quadrant (opening angle $\pi/2$) $\Gamma$ is bounded by $B$ on the boundary and
satisfies $|\Gamma(z)|\le Be^{2\sqrt K|z|}$, i.e.\ it has order at most $1<2 = \pi/(\pi/2)$, so by the Phragm\'en--Lindel\"of principle for a
quadrant, $|\Gamma|\le B$ on that quadrant. Hence $\Gamma$ is bounded on $\mathbb C$, and constant
by Liouville. So $F(s) = F(0)$ for all real $s$ ($s = (\sqrt s)^2$ or $s =
(i\sqrt{-s})^2$), i.e.\ $\ell(\Exp(s\mathfrak D)v) = \ell(v)$. Differentiating at $s=0$,
$\ell(\mathfrak Dv) = 0$. As $\ell$ was arbitrary, $\mathfrak Dv = 0$ by Hahn--Banach.
\prov{Literature: Gelfand 1941, Hille 1944}{A special case of the Gelfand--Hille theorems \cite{Gel41,Hil44} (survey \cite{Zem94}): a bounded group with quasinilpotent generator is trivial; the Phragm\'en--Lindel\"of proof given is self-contained.}

\begin{proof}[Proof of Proposition~\ref{prop:cp}]
$\mathfrak D = [L_x,L_h]$ is a derivation, so $\|\Exp(t\mathfrak D)v\|\le\|v\|$ (Lemma~\ref{lem:expder}), and
$n!\|\mathfrak D^nv\|\le K^n\|v\|$ by VI.6. By VI.7, $\mathfrak D=0$, i.e.\ $x\opc h$.
\end{proof}

\paragraph{VI.8 The chain in a JB-algebra.}
By Lemma~\ref{lem:twosided} and Proposition~\ref{prop:cp}, any $x,h$ with $W_{x,h}$ bounded on
$\mathbb R$ satisfy $x\opc h$; so (SQ) (Lemma~\ref{lem:SQ}) holds in every JB-algebra.

\emph{Powers.} If $a\opc y$ then $a^n\opc y$ for all $n$, by strong induction:
$n = 0,1$ are clear. For $n = 2k$, $k\ge1$: $a^k\opc y$, so (SQ) gives $y\opc(a^k)^2 = a^{2k}$.
For $n = 2k+1$, $k\ge1$: $a^k\opc y$ and $a^{k+1}\opc y$ (as $k+1<n$), so
$a^k+a^{k+1}\opc y$ and (SQ) gives $y\opc(a^k+a^{k+1})^2 = a^{2k} + 2a^{2k+1} + a^{2k+2}$
(power-associativity). Also $y\opc a^{2k}$ and $y\opc a^{2k+2}$ by (SQ) applied to $a^k$
and $a^{k+1}$. Subtracting and dividing by $2$, $y\opc a^{2k+1}$.
\prov{New? (proof); result from their work}{The JB case of Theorem~\ref{thm:main} is \cite[Thm.~3.13]{vdW}, proved via $\cA^{**}$; this bidual-free route we have not found.}

\begin{proof}[Proof of Theorem~\ref{thm:main} for JB-algebras]
By Lemma~\ref{lem:1}, $D(A)=0$; by Proposition~\ref{prop:flow}, $W_{x,A}$ is bounded by
$\|x^2\|$ on $\mathbb R$; by Lemma~\ref{lem:twosided} and Proposition~\ref{prop:cp}, $x\opc A$;
by Lemma~\ref{lem:3}, $y\opc x^2 = a$. By the powers argument $a^n\opc y$ for all $n$, so
$y\in\comm(a)$ (density of polynomials in $C(a)$ and closedness, as in
Proposition~\ref{prop:pows}). Finally, for each $w\in C(a)$ we have $y\opc w$, and (SQ)
gives $w\opc y^2 = b$. So $b\in\comm(a)$.
\end{proof}

\begin{remark}
The same (SQ) argument shows that $\comm(c)$ is closed under the product: if $w\opc y$ and
$w\opc z$ then (SQ) gives $w\opc y^2$, $w\opc z^2$, $w\opc(y+z)^2$, hence
$w\opc 2\,y\jp z$. Together with the theorem this makes $[0,1]_\cA$ with
$a\ast b = U_{\sqrt a}b$ a convex sequential effect algebra, normal when $\cA$ is JBW.
\end{remark}
\prov{Their work: vdW 2020, Prop.~2.34, Thm.~4.2, Prop.~4.4}{$\comm(c)$ a subalgebra is \cite[Prop.~2.34]{vdW} (JB case via $\cA^{**}$); the SEA axioms and normality are \cite[Thm.~4.2, Prop.~4.4]{vdW}. The (SQ) proof of closure is new as far as we know.}

\paragraph{The SEA axioms, one by one.} Let $E = [0,1]_\cA$ with $a\perp b$ iff $a+b\le1$,
$a\oplus b = a+b$, $a^\perp = 1-a$, and $a\ast b = U_{\sqrt a}b$; write $a\mid b$ for
$a\ast b = b\ast a$. The product maps $E\times E$ into $E$, since
$0\le U_{\sqrt a}b\le U_{\sqrt a}1 = a\le1$ by Fact~\ref{f:pos}. The key equivalence is
\[
 (\star)\qquad a\mid b\iff C(b)\subseteq\comm(a)\qquad(a,b\in E):
\]
``$\Rightarrow$'' is the theorem ($b\in\comm(a)$) together with the fact that $\comm(a)$ is a
closed linear subspace containing $1$ (end of \S2) and, by this Remark, a subalgebra;
``$\Leftarrow$'' is the converse in \S1, since $b\in C(b)$. The axioms of Gudder and
Greechie \cite{GG02} (as formalised; (S6) is normality) then hold as follows.
\begin{itemize}
\item[(S1)] $a\ast(b\oplus c) = a\ast b\oplus a\ast c$ for $b\perp c$: $U_{\sqrt a}$ is linear,
 and $a\ast b+a\ast c = a\ast(b+c)\le a\le 1$.
\item[(S2)] $1\ast a = U_1a = a$, since $\sqrt1 = 1$ and $U_1 = I$.
\item[(S3)] $a\ast b = 0\Rightarrow b\ast a = 0$: Fact~\ref{f:S3}.
\item[(S4)] If $a\mid b$ then $a\mid b^\perp$: $1-b\in C(b)\subseteq\comm(a)$ by $(\star)$,
 and the converse in \S1 gives $a\mid 1-b$. Also $(a\ast b)\ast c = a\ast(b\ast c)$: by
 $(\star)$ the elements $a^{1/4},\sqrt a\in C(a)$ and $\sqrt b\in C(b)$ pairwise
 operator-commute, and so do their squares, so $U_{a^{1/4}}$ commutes with $U_{\sqrt b}$.
 By Fact~\ref{f:Uopc}, $a\ast b = a\jp b$ and $\sqrt a\jp\sqrt b = U_{a^{1/4}}\sqrt b\ge0$;
 by the fundamental formula
 $(\sqrt a\jp\sqrt b)^2 = U_{a^{1/4}}U_{\sqrt b}U_{a^{1/4}}1 = U_{\sqrt b}U_{a^{1/4}}\sqrt a
 = U_{\sqrt b}\,a = a\jp b$, so $\sqrt{a\jp b} = \sqrt a\jp\sqrt b$ by uniqueness of square
 roots, and
 $U_{\sqrt{a\jp b}} = U_{a^{1/4}}U_{\sqrt b}U_{a^{1/4}} = U_{\sqrt b}(U_{a^{1/4}})^2 =
 U_{\sqrt b}U_{\sqrt a} = U_{\sqrt a}U_{\sqrt b}$. Hence
 $(a\ast b)\ast c = U_{\sqrt a}U_{\sqrt b}c = a\ast(b\ast c)$.
\item[(S5)] If $c\mid a$ and $c\mid b$ then $c\mid a\ast b$ (for all $a,b$) and, if
 $a\perp b$, $c\mid a\oplus b$: by $(\star)$, $C(a)\cup C(b)\subseteq\comm(c)$, so
 $\sqrt a, b\in\comm(c)$, and since $\comm(c)$ is a subalgebra,
 $a\ast b = 2\sqrt a\jp(\sqrt a\jp b) - a\jp b\in\comm(c)$; also $a+b\in\comm(c)$. The
 converse in \S1 gives $c\mid a\ast b$ and $c\mid a\oplus b$.
\item[(Conv.)] $E$ is convex: $[0,1]_\cA$ is a convex subset of $\cA$ containing $0$, so
 $ta\in E$ for $a\in E$ and $t\in[0,1]$.
\item[(S6)] \emph{Normality, when $\cA$ is JBW.} $E$ is directed complete (monotone
 completeness, \cite[\S4.1]{HOS}). For directed $S\subseteq E$ with supremum $s$,
 $a\ast s = \sup_{t\in S}a\ast t$ by part (iii)(a) of the sketches after
 Fact~\ref{f:jbw} (with $r = \sqrt a\ge0$). If $a\mid t$ for all $t\in S$ then
 $S\subseteq\comm(a)$ by $(\star)$, i.e.\ every $w\in C(a)$ operator-commutes with every
 $t\in S$, hence with $s$ by part (iii)(b); so $s\in\comm(a)$ and $a\mid s$ by the
 converse.
\end{itemize}
\prov{Formalised}{\texttt{jbSEA} and \texttt{sea16\_jb} in \texttt{SEA/JBAllSEA.lean} build
the SEA from the two inputs ``the theorem'' and ``$\comm(c)$ is a subalgebra''
(\texttt{JBCommutation}; S4 via \texttt{Uo\_sqrt\_mul\_of\_compat},
\texttt{sq\_assoc\_of\_compat}; S5 via \texttt{compat\_sq}, \texttt{compat\_add});
\texttt{jbwNormalSEA} adds (S6) from \texttt{jQ\_isLUB} and \texttt{JBWCommSup}. The inputs
are discharged in \texttt{sea16\_jbw\_unconditional} (\texttt{SEA/JBWCommFull.lean}, JBW) and
\texttt{jbCommutation\_all}, \texttt{sea16\_jb\_unconditional\_all}
(\texttt{SEA/JBContPeirce.lean}, all JB-algebras).}

\section*{Appendix A. The fundamental formula and positivity without Macdonald's theorem}
\phantomsection\label{app:mac}
We prove Facts~\ref{f:pos} and~\ref{f:ff}, following \texttt{REC/JBMacdonald.lean}. The
only algebra is one degree-4 identity; the rest is analysis in $B(\cA)$. We use (J),
Facts~\ref{f:norm} and~\ref{f:cfc}, and, in A.5, states: positive unital linear functionals
$\varphi$, which satisfy $|\varphi(v)|\le\|v\|$, and for every $w\ge0$ there is a state with
$\varphi(w) = \|w\|$ (Hahn--Banach in the order-unit space $\cA$). Positivity of $U_a$ is
not used in A.1--A.4.

\paragraph{A.1 The degree-4 identity.} \emph{For all $w,k$:
$2U_{w,k\jp w} = L_kU_w + U_wL_k$.}
Fix $z$. Using commutativity, the two instances of (J) are
\begin{align*}
 j(w,k,w;z) &= 2\,w\jp((k\jp w)\jp z) - 2\,(k\jp w)\jp(w\jp z) + k\jp(w^2\jp z)
   - w^2\jp(k\jp z),\\
 j(w,k,z;w) &= w\jp(w\jp(k\jp z)) - w^2\jp(k\jp z) + k\jp(w\jp(w\jp z))\\
   &\qquad - 2\,(w\jp z)\jp(k\jp w) + z\jp(w\jp(k\jp w)).
\end{align*}
Hence $j(w,k,w;z) - 2j(w,k,z;w)$ equals
\begin{align*}
 &2w\jp((k\jp w)\jp z) + 2(k\jp w)\jp(w\jp z) - 2(w\jp(k\jp w))\jp z\\
 &\qquad - 2k\jp(w\jp(w\jp z)) + k\jp(w^2\jp z) - 2w\jp(w\jp(k\jp z)) + w^2\jp(k\jp z),
\end{align*}
which is $2U_{w,k\jp w}z - k\jp U_wz - U_w(k\jp z)$, since
$k\jp U_wz = 2k\jp(w\jp(w\jp z)) - k\jp(w^2\jp z)$ and
$U_w(k\jp z) = 2w\jp(w\jp(k\jp z)) - w^2\jp(k\jp z)$. By (J) it is $0$.
\prov{Literature}{This is McCrimmon's (FFV) at $y = 1$ \cite[5.2.3]{McC}; the two-instance
derivation is routine. Formalised: \texttt{two\_U2\_mul}.}

\paragraph{A.2 The conjugation formula.} \emph{For all $h,b$:
$U_{\Exp(L_h)b} = \Exp(L_h)\,U_b\,\Exp(L_h)$. In particular $U_{e^h} = \Exp(L_h)^2 =
\Exp(2L_h)$, and $U_{U_{e^h}b} = U_{e^h}U_bU_{e^h}$.}
Put $w(t) = \Exp(tL_h)b$, so $w'(t) = h\jp w(t)$. Since $U_{w+d} = U_w + 2U_{w,d} + U_d$ and
$\|U_d\|\le3\|d\|^2$, the map $t\mapsto U_{w(t)}\in B(\cA)$ is differentiable with
derivative $2U_{w(t),h\jp w(t)}$. Let $F(t) = \Exp(-tL_h)\,U_{w(t)}\,\Exp(-tL_h)$. By the
product rule, and as $L_h$ commutes with $\Exp(-tL_h)$,
\[
 F'(t) = \Exp(-tL_h)\bigl(-L_hU_{w(t)} + 2U_{w(t),h\jp w(t)} - U_{w(t)}L_h\bigr)\Exp(-tL_h)
 = 0
\]
by A.1 with $k = h$. So $F(1) = F(0) = U_b$, and multiplying by $\Exp(L_h)$ on both sides
($\Exp(L_h)\Exp(-L_h) = I$) gives the formula. For $b = 1$: $\Exp(L_h)1 = e^h$
(Fact~\ref{f:cfc}) and $U_1 = I$, so $U_{e^h} = \Exp(L_h)^2 = \Exp(2L_h)$. Finally
$U_{e^h}b = \Exp(L_h)(\Exp(L_h)b)$, so applying the formula twice,
$U_{U_{e^h}b} = \Exp(L_h)^2\,U_b\,\Exp(L_h)^2 = U_{e^h}U_bU_{e^h}$.

\paragraph{A.3 Logarithms.} \emph{If $c\ge\delta1$ with $\delta>0$, then $c = e^g$ for
$g = \log(c)$.} As $(t-\delta)(c) = c-\delta1\ge0$, Fact~\ref{f:cfc} gives
$\spec c\subseteq[\delta,\infty)$, where $\log$ is continuous. Since $f\mapsto f(c)$ is an
isometric algebra homomorphism, $g^n = (\log^n)(c)$, and
$e^g = \sum_n g^n/n! = \bigl(\sum_n\log^n/n!\bigr)(c) = (\exp\circ\log)(c) = c$, the series
of functions converging uniformly on $\spec c$.

\paragraph{A.4 The fundamental formula for all $a$.} Fix $a,b$ and put
$\Phi(\lambda) = U_{U_{a+\lambda1}b} - U_{a+\lambda1}U_bU_{a+\lambda1}\in B(\cA)$. Since
$U_{a+\lambda1} = U_a + 2\lambda L_a + \lambda^2I$, the element
$U_{a+\lambda1}b = U_ab + 2\lambda\,a\jp b + \lambda^2b$ is a polynomial in $\lambda$, and
$\Phi$ is a polynomial in $\lambda$ (of degree $\le4$) with coefficients in $B(\cA)$. For
$\lambda>\|a\|$ we have $a+\lambda1\ge(\lambda-\|a\|)1$ with $\lambda-\|a\|>0$
(Fact~\ref{f:norm}), so $a+\lambda1 = e^g$ for some $g$ by A.3, and $\Phi(\lambda) = 0$ by
A.2. A polynomial that vanishes on an interval has zero coefficients, so $\Phi(0) = 0$:
\[
 U_{U_ab} = U_aU_bU_a ,
\]
and with $b = 1$ ($U_a1 = a^2$, $U_1 = I$), $U_{a^2} = U_a^2$.
\prov{New? (route)}{Analytic proofs of FF exist in finite dimension \cite{FK}, and a
Macdonald-free algebraic proof is van de Wetering's \cite{vdWFF}; we have not found this
Banach-space ODE argument (A.2--A.4) for JB-algebras. The formalisation continues
analytically in $\lambda$ rather than using the polynomial. Formalised:
\texttt{Uo\_exp\_apply}, \texttt{Uo\_eE}, \texttt{ff\_eE}, \texttt{exists\_eE\_eq},
\texttt{Uo\_Uo}, \texttt{Uo\_mul\_self}.}

\paragraph{A.5 Positivity of $\Exp(tL_h)$, and of $U_c$ for $c\ge0$.}
(a) \emph{If $l>0$, $l\|h\|<1$ and $y - l\,h\jp y\ge0$, then $y\ge0$.} Suppose not. Put
$N = \|y\|$ and $w = N1-y\ge0$; $w\ne0$ (else $y = N1\ge0$). Take a state $\varphi$ with
$\varphi(w) = \nu := \|w\|$, and put $z = \nu1-w\ge0$, so $\varphi(z) = 0$ and $y = z+m1$
with $m = N-\nu$; $m<0$, since otherwise $y\ge0$. As $0\le z^2\le\|z\|z$ (functional
calculus: $t^2\le\|z\|t$ on $\spec z\subseteq[0,\|z\|]$), $\varphi(z^2) = 0$. For every
$s\in\mathbb R$, $0\le\varphi((z+sh)^2) = \varphi(z^2) + 2s\,\varphi(h\jp z) +
s^2\varphi(h^2) = 2s\,\varphi(h\jp z)+s^2\varphi(h^2)$, which forces $\varphi(h\jp z) = 0$.
Hence
\[
 \varphi(y - l\,h\jp y) = \varphi(z) + m - l\,\varphi(h\jp z) - l\,m\,\varphi(h)
 = m\bigl(1 - l\,\varphi(h)\bigr) < 0,
\]
as $l\varphi(h)\le l\|h\|<1$. This contradicts $y-l\,h\jp y\ge0$.

(b) \emph{Resolvents.} For $0<l$ with $l\|h\|<1$ we have $\|lL_h\|<1$, so $I-lL_h$ is
invertible (Neumann series), and $(I-lL_h)^{-1}$ is positive: if $x\ge0$ and
$y = (I-lL_h)^{-1}x$, then $y - l\,h\jp y = x\ge0$, so $y\ge0$ by (a).

(c) \emph{$\Exp(tL_h)\ge0$ for all real $t$.} For $t\ge0$,
$\Exp(tL_h) = \lim_n(I-\tfrac tnL_h)^{-n}$ in norm (the exponential formula: with
$\alpha = (I-\frac tnL_h)^{-1}$ and $\beta = \Exp(\frac tnL_h)$ one has
$\|\alpha-\beta\| = o(1/n)$, as $s\mapsto(I-sL_h)^{-1}-\Exp(sL_h)$ has derivative $0$ at
$s = 0$, and $\|\alpha^n-\beta^n\|\le n\|\alpha-\beta\|\max(\|\alpha\|,\|\beta\|)^{n-1}$
with $\max(\|\alpha\|,\|\beta\|)^{n}$ bounded in $n$). By (b) this is a norm limit of
products of positive operators, and the positive operators are closed under products and
norm limits ($\cA_+$ is closed). For $t<0$ use $\Exp(tL_h) = \Exp((-t)L_{-h})$.

(d) \emph{$U_c\ge0$ for $c\ge0$.} If $c\ge\delta1$ with $\delta>0$ then $c = e^g$ (A.3) and
$U_c = \Exp(L_g)\Exp(L_g)\ge0$ by A.2 and (c). For $c\ge0$ and $v\ge0$, and $\varepsilon>0$,
$c+\varepsilon1\ge\varepsilon1$, so $0\le U_{c+\varepsilon1}v = U_cv + 2\varepsilon\,c\jp v +
\varepsilon^2v$; letting $\varepsilon\to0$, $U_cv\ge0$.
\prov{Literature (criterion); New? (minor, its check for $L_h$)}{That positive resolvents
give a positive exponential is Alfsen--Shultz's criterion for order derivations
\cite[(1.82), number unverified]{ASss}; the Cauchy--Schwarz check (a) for $L_h$ we have not
found stated. Formalised: \texttt{nonneg\_of\_resolvent}, \texttt{exp\_mulC\_pos} (with
\texttt{exp\_mem\_posOps\_of\_resolvent} in \texttt{REC/Resolvent.lean}),
\texttt{Uo\_eE\_pos}, \texttt{jQ\_nonneg\_of\_nonneg}.}

\paragraph{A.6 Positivity of $U_a$ for every $a$ (the swap identity).} Let $b\ge0$ and
$c = \sqrt b\ge0$, and put $x = U_a(c^2) = U_ab$ and $y = U_c(a^2)$; $y\ge0$ by A.5(d), as
$a^2\ge0$.

(a) \emph{$x^{n+1} = U_aU_c(y^n)$ for all $n\ge0$.} We use A.4 and $U_v(v^k) = v^{k+2}$
(power-associativity: $2v\jp(v\jp v^k) - v^2\jp v^k = 2v^{k+2}-v^{k+2}$). For $n = 0$:
$U_aU_c1 = U_a(c^2) = x$. For $n = 1$: $x^2 = U_x1 = U_aU_{c^2}U_a1 = U_aU_cU_c(a^2) =
U_aU_c(y)$. For the step $n\to n+2$: $U_x = U_aU_{c^2}U_a = U_aU_cU_cU_a$, so
\begin{align*}
 x^{n+3} = U_x(x^{n+1}) &= U_aU_cU_cU_aU_aU_c(y^n) = U_aU_c\,(U_cU_{a^2}U_c)(y^n)\\
 &= U_aU_c\,U_{U_c(a^2)}(y^n) = U_aU_c(y^{n+2}).
\end{align*}
(b) By linearity, $x\jp q(x) = U_aU_c(q(y))$ for every real polynomial $q$.

(c) Let $f(t) = \max(-t,0)$, $R = \max(\|x\|,\|y\|)$ and $\varepsilon>0$, and take a
polynomial $q$ with $|q-f|\le\varepsilon$ on $[-R,R]$ (Weierstrass). As $f = 0$ on
$\spec y\subseteq[0,R]$, $\|q(y)\|\le\varepsilon$; and $\|q(x)-f(x)\|\le\varepsilon$. So
$x\jp f(x) = U_aU_c(q(y)) - x\jp(q(x)-f(x))$ has norm
$\le(\|U_aU_c\|+\|x\|)\varepsilon$. Hence $x\jp f(x) = 0$, i.e.\ the function
$t\max(-t,0)$ vanishes on $\spec x$, so $\spec x\subseteq[0,\infty)$ and $U_ab = x\ge0$.
\prov{New? (in this form)}{In $B(H)$, with $T = ac$, (a) is
$(TT^*)^{n+1} = T(T^*T)^nT^*$; we have not found this Jordan form of the argument.
Formalised: \texttt{swap\_pow}, \texttt{swap\_poly}, \texttt{jQ\_nonneg}, \texttt{hos336}.}

\section*{Verification and formalisation}

\paragraph{Numerical and symbolic checks.}
The identities of Part I were checked exactly in the free commutative non-associative
algebra over $\mathbb Q$: the derivation property equals $j(v,x,u;y) - j(u,y,v;x)$, Lemma~\ref{lem:1}
equals $\tfrac12j(y,x,x;y)-\tfrac12j(y,y,x;x)$, Lemma~\ref{lem:3} equals $j(x,x,y;z)$, and
Lemma~\ref{lem:2} equals both the seven-instance combination above and the ten-instance combination
of the formal proof; and the degree-4 identity A.1 equals $j(w,k,w;z) - 2j(w,k,z;w)$.
Numerically (random real symmetric $5\times5$ matrices, and random
elements of the Albert algebra $H_3(\mathbb O)$), the following vanish to rounding error
(maximal residual $3\cdot10^{-13}$, resp.\ $5\cdot10^{-11}$): (J), the derivation property, Lemmas~\ref{lem:1}, \ref{lem:3} and~\ref{lem:2},
$[L_x,L_{x^2}]$, the instance $\cR$ of (II), $J_3$, the pre-final form
$\cK = -2[X,[X,\cP]] + [X,[\cQ,Y]]$, the fundamental formula, $U_{x^2} = U_x^2$,
$2U_{w,k\jp w} = L_kU_w + U_wL_k$, the square identity in Fact~\ref{f:S3}, the cross term
VI.3 (with $a\jp b = 0$), (Mb) on $C(h)$, and $\Delta = \tfrac12(U_{1,g^2}-U_g)$.
For the additions (Appendix~A, the sketches after Fact~\ref{f:jbw}, and the SEA
paragraph) the following were checked in both algebras (maximal residual $2\cdot10^{-14}$,
resp.\ $2\cdot10^{-11}$): $U_{w+d} = U_w+2U_{w,d}+U_d$, $U_{\Exp(L_h)b} =
\Exp(L_h)U_b\Exp(L_h)$, $U_{e^h} = \Exp(2L_h)$, the fundamental formula at $e^h$, the swap
identity $x^{n+1} = U_aU_c(y^n)$ for $n\le4$ (relative residual), $NL_v =
U_{(N+v)/2}-U_{(N-v)/2}$, $U_{r+\varepsilon1}-U_r = 2\varepsilon L_r+\varepsilon^2I$,
$U_{1+tc} = I+2tL_c+t^2U_c$, and, for $a = e^h$, $b = e^{2h}$: $\sqrt a\jp\sqrt b =
U_{a^{1/4}}\sqrt b$, $(\sqrt a\jp\sqrt b)^2 = a\jp b$ and $U_{\sqrt{a\jp b}} =
U_{\sqrt a}U_{\sqrt b}$. In $\mathrm{Sym}_5(\mathbb R)$ we also checked $e^{\log c} = c$, the
positivity of $\Exp(\pm1.5L_h)b$ and $(I-lL_h)^{-1}b$ for $b\ge0$, and the convergence of
the exponential formula; and, pointwise on a grid, the scalar inequalities behind (i) and
(ii) after Fact~\ref{f:jbw}.

\paragraph{Formalised as.}
{\raggedright The proof is machine-checked in Lean. Part I: \texttt{jb\_lin}, \texttt{der\_apply},
\texttt{der\_xy}, \texttt{jk\_apply}/\texttt{jk}, \texttt{opComm\_sq\_of\_opComm\_mul}.
Part II: \texttt{exp\_der\_mul}, \texttt{exp\_der\_norm}, \texttt{flow},
\texttt{flow\_bound}. Part III: \texttt{gap\_zero}, \texttt{ph\_sproj\_of\_flow},
\texttt{opComm\_of\_sprojs}, \texttt{pow\_opComm\_of\_sprojs}, \texttt{mem\_comm\_of\_pows}.
Parts IV--V: \texttt{opComm\_sqrt\_of\_sq}, \texttt{opComm\_pow\_of\_opComm},
\texttt{jbw\_mul\_mem}, \texttt{jbw\_mem\_comm}. Part VI: \texttt{gap\_zero\_s},
\texttt{gaps\_of\_flow\_bdd}, \texttt{Mb}, \texttt{UU\_zero}, \texttt{UL\_zero},
\texttt{mul\_eq\_zero\_of\_Uo}, \texttt{ff\_cross}, \texttt{u2\_sep},
\texttt{piece\_bound}, \texttt{delOp\_eq\_zero}, \texttt{ks\_norm},
\texttt{eq\_zero\_of\_bounded\_group}, \texttt{contPeirce}, \texttt{opComm\_sq},
\texttt{opComm\_pow}, \texttt{jb\_mem\_comm}, \texttt{jb\_mul\_mem}.
Standard facts: \texttt{JBCalculus}, \texttt{JBMacdonald} (\texttt{two\_U2\_mul},
\texttt{Uo\_exp\_apply}, \texttt{exists\_eE\_eq}, \texttt{Uo\_Uo}, \texttt{Uo\_eE},
\texttt{nonneg\_of\_resolvent}, \texttt{exp\_mulC\_pos}, \texttt{swap\_pow},
\texttt{jQ\_nonneg}), \texttt{JBPeirce}, \texttt{JBWProj} (\texttt{sproj},
\texttt{spectral\_approx\_norm}), \texttt{opComm\_of\_isLUB},
\texttt{sq\_eq\_zero\_comm}, \texttt{ph\_of\_jQ\_one\_sub}, \texttt{jQ\_isLUB}.\par}

\begin{thebibliography}{99}
\small
\bibitem{HOS} H.~Hanche-Olsen and E.~St{\o}rmer, \emph{Jordan Operator Algebras},
Pitman, 1984.
\bibitem{vdW} J.~van de Wetering, \emph{Commutativity in Jordan operator algebras},
J.\ Pure Appl.\ Algebra 224(11) (2020) 106407; arXiv:1912.01903. (Numbering as in
arXiv v1, the only version.)
\bibitem{vdW19} J.~van de Wetering, \emph{Sequential product spaces are Jordan algebras},
J.\ Math.\ Phys.\ 60 (2019) 062201; arXiv:1803.11139.
\bibitem{vdWFF} J.~van de Wetering, \emph{An algebraic semi-automated proof of the
fundamental identity of Jordan algebras}, arXiv:1807.00164, 2018.
\bibitem{WWvdW20} A.~Westerbaan, B.~Westerbaan and J.~van de Wetering, \emph{The three
types of normal sequential effect algebras}, Quantum 4 (2020) 378; arXiv:2004.12749.
\bibitem{WWvdW21} B.~Westerbaan and J.~van de Wetering, \emph{A computer scientist's
reconstruction of quantum theory}, arXiv:2109.10707.
\bibitem{AS} E.~M.~Alfsen and F.~W.~Shultz, \emph{Geometry of State Spaces of Operator
Algebras}, Birkh\"auser, 2003.
\bibitem{ASss} E.~M.~Alfsen and F.~W.~Shultz, \emph{State Spaces of Operator Algebras:
Basic Theory, Orientations, and C$^*$-products}, Birkh\"auser, 2001.
\bibitem{FK} J.~Faraut and A.~Kor\'anyi, \emph{Analysis on Symmetric Cones}, Oxford
University Press, 1994.
\bibitem{McC} K.~McCrimmon, \emph{A Taste of Jordan Algebras}, Springer, 2004.
\bibitem{McC66} K.~McCrimmon, \emph{A general theory of Jordan rings}, Proc.\ Nat.\ Acad.\
Sci.\ USA 56 (1966) 1072--1079.
\bibitem{Jac49} N.~Jacobson, \emph{Derivation algebras and multiplication algebras of
semi-simple Jordan algebras}, Ann.\ of Math.\ 50 (1949) 866--874.
\bibitem{Jac68} N.~Jacobson, \emph{Structure and Representations of Jordan Algebras},
AMS Colloquium Publ.\ 39, 1968.
\bibitem{Jac69} N.~Jacobson, \emph{Lectures on Quadratic Jordan Algebras}, Tata Institute,
Bombay, 1969.
\bibitem{Jac89} N.~Jacobson, \emph{Operator commutativity in Jordan algebras}, in:
\emph{Collected Mathematical Papers}, Springer, 1989, 169--172.
\bibitem{Mac60} I.~G.~Macdonald, \emph{Jordan algebras with three generators}, Proc.\
London Math.\ Soc.\ (3) 10 (1960) 395--408.
\bibitem{Shu79} F.~W.~Shultz, \emph{On normed Jordan algebras which are Banach dual
spaces}, J.\ Funct.\ Anal.\ 31 (1979) 360--376.
\bibitem{GN01} S.~Gudder and G.~Nagy, \emph{Sequential quantum measurements}, J.\ Math.\
Phys.\ 42 (2001) 5212--5222.
\bibitem{GG02} S.~Gudder and R.~Greechie, \emph{Sequential products on effect algebras},
Rep.\ Math.\ Phys.\ 49 (2002) 87--111.
\bibitem{Fug50} B.~Fuglede, \emph{A commutativity theorem for normal operators}, Proc.\
Nat.\ Acad.\ Sci.\ USA 36 (1950) 35--40.
\bibitem{Put51} C.~R.~Putnam, \emph{On normal operators in Hilbert space}, Amer.\ J.\
Math.\ 73 (1951) 357--362.
\bibitem{Ros58} M.~Rosenblum, \emph{On a theorem of Fuglede and Putnam}, J.\ London Math.\
Soc.\ 33 (1958) 376--377.
\bibitem{Kle57} D.~C.~Kleinecke, \emph{On operator commutators}, Proc.\ Amer.\ Math.\
Soc.\ 8 (1957) 535--536.
\bibitem{Shi56} F.~V.~Shirokov, \emph{Proof of a conjecture of Kaplansky}, Uspekhi Mat.\
Nauk 11(4) (1956) 167--168.
\bibitem{Gel41} I.~Gelfand, \emph{Zur Theorie der Charaktere der Abelschen topologischen
Gruppen}, Mat.\ Sb.\ 9(51) (1941) 49--50.
\bibitem{Hil44} E.~Hille, \emph{On the theory of characters of groups and semi-groups in
normed vector rings}, Proc.\ Nat.\ Acad.\ Sci.\ USA 30 (1944) 58--60.
\bibitem{Zem94} J.~Zem\'anek, \emph{On the Gelfand--Hille theorems}, in: Functional
Analysis and Operator Theory, Banach Center Publ.\ 30 (1994) 369--385.
\end{thebibliography}

\end{document}
